\documentclass[preprint,12pt]{elsarticle}

\usepackage[T1]{fontenc}
\usepackage[utf8]{inputenc}
\usepackage{amsmath}
\usepackage{amssymb}
\usepackage{graphicx}
\usepackage{subcaption}
\usepackage{booktabs}
\usepackage{url}
\usepackage{mathtools}

\usepackage{caption}

\usepackage{array}

\usepackage{xcolor}

\journal{Information and Software Technology}

\begin{document}

\begin{frontmatter}

\title{Property-Based Testing of ML-Based Anomaly Detectors under Incomplete Ground Truth}

\author[inst1]{Manuel Muñoz}
\ead{mmr411@ual.es}

\author[inst1]{Antonio Becerra-Terón\corref{cor1}}
\ead{abecerra@ual.es}

\author[inst1]{Manuel Torres}
\ead{mtorres@ual.es}

\author[inst1]{Francisco Rodríguez}
\ead{frrodrig@ual.es}

\cortext[cor1]{Corresponding author}

\address[inst1]{Department of Informatics, CIESOL-ceiA3,
University of Almería, Carretera de Sacramento s/n,
04120 Almería, Spain}



\begin{abstract}
Machine-learning (ML)-based anomaly detectors are difficult to test in
real-world Internet of Things (IoT) systems because exhaustive ground-truth
labels for heterogeneous anomaly patterns are rarely available. This paper
investigates property-based testing (PBT) as a complementary testing strategy
in which domain-grounded anomaly knowledge is transformed into executable
properties, property-specific generators, and partial behavioural oracles.
Rather than constructing a fully labelled synthetic test dataset, the proposed
framework generates test cases from property-specific generation spaces and
checks explicit behavioural expectations of the detector. The framework is instantiated for seven anomaly properties in greenhouse
climatic sensor data and evaluated against a frozen LightGBM-based anomaly
detection and classification pipeline serving as the system under test, with
strong conventional predictive performance. The resulting campaign executes
26,767 property checks over real admissible base contexts, preserving
case-level traceability from generation parameters to system output and oracle
outcome. The evaluated system conforms to 24,231
checks (90.53\%) and produces 2,536 counterexamples, with strong
property-dependent variation: missing values, out-of-range values, and partial
system failures show complete conformance, whereas stuck sensors and
correlation deviations expose substantial non-conformance. Stratum-level
analysis further shows that counterexamples are concentrated in specific
regions of the evaluated generation spaces rather than uniformly distributed.
These results show that PBT provides structured behavioural evidence for
ML-based anomaly detectors under incomplete ground truth, complementing
aggregate predictive metrics with traceable conformance and counterexample
information.
\end{abstract}

\begin{keyword}
Property-based testing \(\sep\) Machine learning testing \(\sep\) Anomaly detection \(\sep\) Test oracle \(\sep\) Incomplete ground truth
\end{keyword}

\end{frontmatter}


\section{Introduction}
\label{sec:introduction}

Testing machine-learning-based systems poses distinctive software
engineering challenges, particularly regarding test-input generation,
behavioural adequacy, and the definition of reliable test oracles
\cite{riccio2020testing,benbraiek2020testingml,zhang2022mltesting}.
These challenges are part of the broader problem of establishing
trustworthy and verifiable behaviour in AI-enabled systems
\cite{seshia2022toward,huang2020safety}. In particular, determining whether
the behaviour produced by a learned model is correct is closely related to
the long-standing test oracle problem in software testing, which arises
when the expected output for an arbitrary input is difficult or costly to
determine \cite{barr2015oracle}. This difficulty is especially relevant for
machine-learning-based systems, whose behaviour is learned from data rather
than completely specified through explicit program logic.

Conventional predictive evaluation mitigates this difficulty when labelled
observations are available by comparing model outputs against known outcomes
and summarizing their agreement through metrics such as precision, recall,
F1-score, or area-under-the-curve measures. Such evaluation provides
valuable evidence when representative ground-truth labels are available,
but its scope is constrained when the expected behaviour of the system
cannot be exhaustively specified over its operational input space
\cite{riccio2020testing}.

This limitation is particularly relevant to anomaly detection
\cite{adhikari2024iotAnomalySurvey,gkountakos2024comprehensive}.
Real-world anomaly-detection datasets often contain abundant observations
of normal operation but comparatively limited ground-truth information
describing the heterogeneous ways in which abnormal behaviour may occur.
The problem is particularly challenging in sensor-based systems, where
anomalies are not necessarily pointwise deviations: they may emerge from
temporal persistence, physical constraints, operating context, or
relationships among multiple variables
\cite{karkouch2016data,teh2020sensor,zou2023sensorfault}. Moreover, rare or
extreme observations are not necessarily erroneous, while apparently
admissible measurements may constitute anomalies when interpreted in their
temporal or physical context. Consequently, exhaustive manual labeling of
the anomaly space is generally impractical, limiting the behaviours that
can be assessed through conventional predictive evaluation.

When sufficiently representative anomaly labels are unavailable, controlled
anomalous observations can be constructed and used as synthetic ground
truth. This makes it possible to evaluate learned anomaly detectors against
known anomaly classes using conventional predictive metrics. We previously
followed this strategy in a taxonomy-driven framework for greenhouse
climatic sensor data, where empirically characterized anomaly patterns were
used to generate controlled anomalous observations for training and
evaluating supervised detectors \cite{carrero2026ifac}. This strategy
provides a basis for predictive evaluation, but it also raises a broader
software-testing question: whether the same domain knowledge can be
operationalized directly as executable test specifications, enabling
systematic exploration of anomaly behaviours without first materializing
them as a fixed labelled test dataset.

Property-based testing (PBT) provides a natural mechanism for investigating
this question. Rather than expressing expected behaviour exclusively through
individual input--output examples, PBT defines general properties that are
exercised over automatically generated test cases
\cite{claessen2000quickcheck,almendros2017automatic}. Constrained generators
explore relevant families of inputs, while failed property checks produce
concrete counterexamples for further investigation
\cite{almendros2017automatic,goldstein2024pbt}. This perspective suggests an
alternative testing cycle in which domain knowledge directly supports
executable test specifications for machine-learning-based anomaly detectors.

Under this perspective, anomaly specifications do not merely support the
construction of labelled observations; they define constrained input spaces
and behavioural expectations that can be exercised through property-specific
generators and partial behavioural oracles. This does not eliminate the test
oracle problem or replace conventional predictive evaluation. Instead, it
enables selected behaviours to be tested under explicitly defined conditions
without requiring a complete oracle over the operational input space. Related
approaches, particularly metamorphic testing, similarly exploit known
relations or domain properties when complete test oracles are unavailable
\cite{segura2016metamorphic,riccio2020testing}. PBT provides a complementary
mechanism in which executable properties directly guide test generation and
counterexample discovery.

We investigate this testing strategy in the context of anomaly detection
for IoT climatic sensor data. Our previous work established an empirically
grounded taxonomy comprising missing values, stuck sensors, noise,
out-of-range values, contextual anomalies, correlation deviations, and
valid extreme values \cite{carrero2026ifac}. A subsequent predicate-based
reliability model formalized these anomaly semantics over sensor
observations, temporal windows, and contextual relationships to support
anomaly-specific correction and reliability assessment
\cite{carrero2026qrs}. Together, these studies provide the domain knowledge
underlying the present work. Rather than using anomaly semantics to construct
labelled observations or support correction, we investigate their use as
executable specifications for systematically testing a learned anomaly
detector.

Building on this foundation, we propose a property-based testing framework
that transforms domain-grounded anomaly patterns into executable test
specifications combining behavioural expectations, property-specific
generators, and partial behavioural oracles. The generators systematically
explore constrained input spaces by varying semantically meaningful
dimensions of each anomaly pattern, such as the affected sensor, magnitude,
duration, or operational context, while property-specific conditions restrict
generation to cases for which the expected behaviour can be meaningfully
evaluated. Each valid test case is executed against the system under test
(SUT), and violations of the expected behaviour are retained as
counterexamples. Beyond aggregate conformance, this process enables the
characterization of how SUT behaviour varies across predefined dimensions of
the corresponding generation spaces, revealing where the expected behaviour
is preserved or violated.

The framework is evaluated using real-world greenhouse IoT sensor data and
an anomaly detection and classification system previously shown to achieve
strong predictive performance under conventional ground-truth-based
evaluation. This provides a particularly informative empirical setting for
investigating whether PBT can reveal behavioural evidence beyond aggregate
predictive metrics, including counterexamples and variations in conformance
across property-defined input spaces. The study therefore examines PBT as a
complementary testing perspective rather than a replacement for conventional
predictive evaluation.

The main contributions of this work are:
\begin{itemize}

    \item A property-based testing framework that transforms domain-grounded
    anomaly knowledge into executable test specifications for anomaly
    detection systems.

    \item Property-specific generators and partial behavioural oracles that
    systematically exercise semantically constrained anomaly spaces without
    requiring complete ground truth over the operational input space.

    \item An empirical evaluation on an anomaly detection system with
    previously demonstrated strong predictive performance, characterizing
    behavioural conformance and localizing counterexamples across semantically
    defined regions of the property-specific generation spaces.

\end{itemize}

The remainder of this paper is organized as follows.
Section~\ref{sec:related_work} reviews the relevant literature, and
Section~\ref{sec:framework} presents the proposed framework.
Section~\ref{sec:empirical_study} describes the empirical study and research
questions, while Section~\ref{sec:results} reports the results.
Section~\ref{sec:discussion} discusses their implications, generalizability,
and threats to validity. Finally, Section~\ref{sec:conclusions} concludes the
paper and outlines future work.


\section{Background and Related Work}
\label{sec:related_work}

\subsection{Testing Machine Learning Systems}

Testing ML-based systems differs from testing conventional software because
observed behaviour is determined not only by source code, but also by training
data, learned decision logic, preprocessing, and the operational data
distribution. Systematic mappings and surveys identify test-input generation,
oracle construction, data adequacy, robustness, and behavioural specification
as recurring challenges in ML testing
\cite{riccio2020testing,zhang2022mltesting,benbraiek2020testingml}.

Conventional model evaluation provides essential evidence about predictive
performance over labelled data or validation distributions. However, aggregate
metrics such as accuracy, precision, recall, or F1-score do not directly
characterize system behaviour under specific operational conditions or
semantically meaningful input variations. Early systematic testing approaches
extend evaluation beyond aggregate predictive metrics. DeepXplore combines
neuron-coverage-guided input generation with differential testing to expose
failure-inducing inputs through inconsistent predictions across multiple deep
neural networks \cite{pei2017deepxplore}. DeepTest systematically generates
transformed driving scenes under controlled environmental conditions and uses
metamorphic relations to identify failure-inducing behaviours of
autonomous-driving deep neural networks across different transformation
conditions \cite{tian2018deeptest}. Behavioural testing approaches such as
CheckList organize tests around explicit linguistic capabilities and expected
behaviours, using controlled test construction and perturbation to evaluate
model failures across capabilities and test types
\cite{ribeiro2020checklist}. Recent empirical studies of ML testing practice
similarly report the use of conventional tests, data checks, task-specific
validation code, and domain-specific expectations in ML-enabled systems
\cite{openja2024mlTestingWild,cannavale2026mlTestingPractice}.

From a software-testing perspective, the resulting object of evaluation is
therefore not only the learned model, but the system in which data processing,
learned decision logic, and surrounding software jointly determine observable
behaviour. Testing such systems requires both suitable test inputs and
expectations against which their outputs can be assessed. This leads directly
to the test-oracle problem, particularly when complete ground truth is
unavailable or only partially specified.


\subsection{Testing under Incomplete Ground Truth}

The test-oracle problem arises when the expected output for a test input is
difficult, costly, or impossible to determine \cite{barr2015oracle}. In
ML-based systems, the problem is especially relevant because the correct
response may depend on uncertain labels, statistical behaviour, context, or
domain-specific interpretation. A complete oracle would determine the
correctness of every SUT output, but many practical testing settings only
support weaker forms of evidence. Existing strategies include pseudo-oracles,
differential testing, learned oracles, specification-based checks, and
metamorphic testing \cite{fontes2021oracles,segura2016metamorphic}.

These mechanisms are related but not equivalent. Differential testing compares
outputs produced by multiple implementations or models, while pseudo-oracles
use an independently constructed mechanism as a surrogate expected-output
source. Metamorphic testing instead exploits expected relations between
executions: rather than requiring the exact output for each individual input,
it checks whether source and follow-up executions satisfy a metamorphic
relation \cite{CHEN20031,segura2016metamorphic}.

For ML classifiers, controlled follow-up test cases can be derived from source
cases through predefined metamorphic relations, with violations providing
concrete failure evidence that can be compared across relations and
classifiers \cite{xie2011metamorphicML}. METTLE extends metamorphic testing to
unsupervised ML systems by defining metamorphic relations and controlled
follow-up inputs whose satisfaction or violation provides evidence about
system behaviour without requiring complete labelled ground truth
\cite{xie2020mettle}. Its evaluation compares metamorphic-relation behaviour
across the tested systems, providing structured failure evidence across
predefined relations.

For the testing setting considered here, partial behavioural oracles provide
a different form of evidence. They evaluate whether an observable SUT
response conforms to a specified expectation under defined conditions; they
do not establish universal correctness of the prediction. This distinction
is important under incomplete ground truth: partial expectations can make
selected behaviours testable, but they must be explicitly scoped and
justified. The practical question is therefore how such expectations can be
made executable and systematically exercised over relevant input spaces.


\subsection{Property-Based Testing for ML Systems}

PBT addresses testing through executable properties rather than only through
manually selected input--output examples. The approach was established by
QuickCheck, which introduced the combination of executable properties,
automated test generation, and counterexample discovery
\cite{claessen2000quickcheck}. Modern PBT frameworks and studies emphasize
that its effectiveness depends not only on automated input generation, but
also on the quality of the properties and the adequacy of the generators used
to explore the input space
\cite{maciver2019hypothesis,lampropoulos2019coverageGuidedPBT,goldstein2024pbt}.
Counterexamples provide concrete evidence of property violations and can be
retained for diagnosis, regression testing, or specification refinement.

For ML-based systems, PBT should not be reduced to randomized input
generation. Properties must express behavioural expectations that are
meaningful for the SUT, while generators must produce inputs for which those
expectations are valid. Recent empirical studies of PBT in practice and in ML
projects highlight both its potential and the difficulty of designing
effective properties and data generators
\cite{wauters2024pbtMLProjects,oliveira2026pbtPython}. Work on generator
coverage and verification further reinforces that the generation mechanism
itself is a critical part of the testing process
\cite{zhou2023coverageGenerators}.

Applying PBT to data-intensive and ML-based systems requires properties that
go beyond the syntactic interface of the model. Meaningful properties often
depend on domain knowledge, empirical evidence, or behavioural requirements
that determine which inputs should be generated and what response should be
expected. This requirement becomes particularly relevant in anomaly detection,
where the same observation may be normal or anomalous depending on temporal
persistence, physical bounds, contextual conditions, or relationships among
variables.


\subsection{Testing ML-Based Anomaly Detection}

Evaluation of anomaly detectors commonly relies on labelled or injected
anomalies, benchmark datasets, and aggregate predictive metrics. General
anomaly-detection surveys and time-series evaluations show that assessment
depends strongly on the definition of anomaly types, label quality, temporal
structure, and metric choice
\cite{10.1145/1541880.1541882,schmidl2022timeeval}. Benchmark critiques also
show that ground-truth quality and dataset artefacts can substantially affect
conclusions about anomaly-detection progress
\cite{wu2021benchmarks}. These concerns are particularly relevant for sensor
systems, where data-quality problems include missingness, physical
implausibility, temporal persistence, and contextual inconsistency
\cite{karkouch2016data,teh2020sensor}.

Controlled anomaly injection and synthetic anomaly generation can provide
known anomalous cases for evaluation, but they do not automatically define a
software-testing oracle. A generated anomaly may support predictive evaluation
if treated as labelled data, whereas behavioural testing additionally requires
explicit conditions under which the generated case is valid and an expected
SUT response against which the output can be checked. AutoConf uses
metamorphic relations and systematically generated configurations to assess
configuration-dependent behaviour without requiring a complete expected-output
oracle \cite{shar2023autoconf}. Metamorphic violations provide structured
evidence of problematic configurations.

Our previous predicate-based reliability work formalized anomaly semantics as
executable predicates for correction and reliability assessment in IoT
climatic sensor data \cite{carrero2026qrs}. The present study builds on this
predicate-oriented foundation from a software-testing perspective.


\subsection{Research Gap}

Prior work provides substantial foundations for testing ML-based systems
under incomplete or unavailable ground truth, including differential,
metamorphic, behavioural, and property-based testing approaches.
Table~\ref{tab:related_work_capability_matrix} compares representative
approaches closest to the capabilities considered in this study.

\begin{table*}[t]
\centering
\caption{Capability-based comparison of related approaches for testing
ML-based systems under incomplete ground truth.
\(\checkmark\) = explicitly supported; \(\circ\) = supported in a related
but not equivalent form; -- = not provided. Abbreviations: spec. =
specification, gen. = controlled generation, CE = counterexamples, eval. =
evaluation, conf. = conformance, char. = characterization, and W-space =
within-space.}
\label{tab:related_work_capability_matrix}
\scriptsize
\setlength{\tabcolsep}{2.6pt}
\renewcommand{\arraystretch}{1.15}
\begin{tabular}{
p{0.22\textwidth}
>{\centering\arraybackslash}p{0.10\textwidth}
>{\centering\arraybackslash}p{0.09\textwidth}
>{\centering\arraybackslash}p{0.08\textwidth}
>{\centering\arraybackslash}p{0.08\textwidth}
>{\centering\arraybackslash}p{0.10\textwidth}
>{\centering\arraybackslash}p{0.12\textwidth}
}
\toprule
\textbf{Work / approach} &
\shortstack{\textbf{Domain}\\\textbf{spec.}} &
\textbf{Gen.} &
\shortstack{\textbf{Partial}\\\textbf{oracle}} &
\textbf{CE} &
\shortstack{\textbf{Conf.}\\\textbf{eval.}} &
\shortstack{\textbf{W-space}\\\textbf{conf.}\\\textbf{char.}} \\
\midrule

DeepXplore~\cite{pei2017deepxplore}
& -- & \(\circ\) & \(\circ\) & \(\checkmark\) & \(\circ\) & -- \\

DeepTest~\cite{tian2018deeptest}
& \(\circ\) & \(\checkmark\) & \(\circ\) & \(\checkmark\) & \(\circ\) & \(\circ\) \\

CheckList~\cite{ribeiro2020checklist}
& \(\checkmark\) & \(\checkmark\) & \(\checkmark\) & \(\checkmark\) &
\(\checkmark\) & \(\circ\) \\

Metamorphic testing for ML classifiers~\cite{xie2011metamorphicML}
& \(\circ\) & \(\checkmark\) & \(\circ\) & \(\checkmark\) &
\(\checkmark\) & \(\circ\) \\

METTLE~\cite{xie2020mettle}
& \(\circ\) & \(\checkmark\) & \(\circ\) & \(\circ\) &
\(\checkmark\) & \(\circ\) \\

AutoConf~\cite{shar2023autoconf}
& \(\circ\) & \(\checkmark\) & \(\circ\) & \(\circ\) &
\(\circ\) & -- \\

\textbf{Our work}
& \(\checkmark\) & \(\checkmark\) & \(\checkmark\) & \(\checkmark\) &
\(\checkmark\) & \(\checkmark\) \\

\bottomrule
\end{tabular}
\end{table*}

The comparison shows that prior work already supports many of the individual
capabilities relevant to this study. In particular, existing approaches
provide mechanisms for controlled test generation, behavioural expectations,
oracle support under incomplete ground truth, failure identification, and
conformance-related evaluation. Several approaches also examine behavioural
variation across structured test conditions, although not as systematic
conformance characterization within property-specific generation spaces.

However, the literature provides limited evidence of an integrated PBT
workflow for ML-based anomaly detectors in which domain-grounded anomaly
semantics jointly determine executable properties, property-specific
controlled generation, and partial behavioural oracles. Such a workflow also
connects generated test cases and counterexamples to conformance evaluation
and to the characterization of behavioural variation across predefined
dimensions of the corresponding generation spaces.

The gap is therefore not the absence of test generation, metamorphic
relations, behavioural expectations, or failure discovery in prior work, but
the limited evidence on how these capabilities can be integrated into a
traceable PBT process for ML-based anomaly detection when expected behaviour
is local, domain-specific, and only partially observable. This is the
methodological gap addressed by the present study.

\section{Property-Based Testing Framework for ML-Based Anomaly Detectors}
\label{sec:framework}

This section presents a property-based testing framework for ML-based anomaly
detectors when complete instance-level ground truth is unavailable. The
framework uses partial behavioural knowledge to define executable properties,
generate property-driven test cases, evaluate SUT behaviour through
property-specific partial oracles, and identify counterexamples. The following
subsections formalize these components and their relationships.

\subsection{Problem Formulation}
\label{sec:problem_formulation}

The ML-based anomaly detector is considered the system under test (SUT).
Let \(D\) denote such a detector, represented as
\begin{equation}
    D : \mathcal{X} \rightarrow \mathcal{Y},
    \label{eq:sut}
\end{equation}
where \(\mathcal{X}\) and \(\mathcal{Y}\) denote the operational input and
output spaces, respectively. For any \(x \in \mathcal{X}\),
\(D(x) \in \mathcal{Y}\) denotes the observable output produced by the SUT.
Depending on the detector, \(x\) may represent an individual or multivariate
observation, a temporal window, a sequence, or an input representation
incorporating contextual information, while \(D(x)\) may represent a
detection decision, an anomaly category, a score, or a structured output.

If complete instance-level ground truth were available, an ideal test oracle
could determine the correctness of the SUT output for any input. Such a
complete oracle can be represented as
\begin{equation}
    O :
    \mathcal{X} \times \mathcal{Y}
    \rightarrow
    \{\mathsf{true},\mathsf{false}\},
    \label{eq:complete_oracle}
\end{equation}
where \(O(x,D(x))\) indicates whether the output produced by the detector for
input \(x\) is correct. The testing problem addressed here arises because
\(O(x,D(x))\) cannot be evaluated for arbitrary operational inputs when
complete instance-level ground truth is unavailable.

Nevertheless, partial domain knowledge may still support conditional
expectations about specific aspects of detector behaviour. Let \(p\)
denote such a behavioural specification and \(\Phi_p(x)\) the conditions
under which its expectation applies. Let \(B_p(x,D(x))\) be a Boolean predicate indicating whether the
observable output \(D(x)\) produced by the SUT for input \(x\) satisfies the
behavioural expectation specified by \(p\). A partial behavioural
specification is therefore expressed as
\begin{equation}
    \Phi_p(x)
    \Rightarrow
    B_p\bigl(x,D(x)\bigr).
    \label{eq:behavioural_specification}
\end{equation}
Unlike \(O\), this specification does not determine the complete correctness
of \(D(x)\), but only expresses a conditional expectation about a specific
aspect of detector behaviour. Such an expectation can therefore support
a property-specific partial oracle, denoted by \(O_p\), rather than a complete
correctness judgement. To enable systematic testing, however, the partial
behavioural specification must first be transformed into an executable property.

\subsection{From Behavioural Specifications to Executable Properties}
\label{sec:executable_properties}

An executable property requires two elements: the conditions determining
\emph{when} the property can be meaningfully tested and the behavioural
expectation specifying \emph{what} is expected from the SUT.
In the behavioural specification introduced in
Eq.~\ref{eq:behavioural_specification}, these elements are represented by
\(\Phi_p(x)\) and \(B_p(x,D(x))\), respectively.

While \(B_p(x,D(x))\) directly expresses the behavioural expectation to be
evaluated for a concrete SUT execution, the conditions represented by
\(\Phi_p(x)\) must first be operationalized to determine when that expectation
can be meaningfully tested. For ML-based anomaly detectors, this involves determining whether an input
represents the target detection situation, constitutes a valid test case, and
supports an unambiguous behavioural expectation.

Accordingly, three complementary conditions are considered. The applicability
condition \(C_p(x)\) determines whether the input exhibits the target detection
situation. The admissibility condition \(A_p(x)\) determines whether the
additional constraints required for a valid test case are satisfied. The
exclusion condition \(E_p(x)\) identifies situations in which alternative
interpretations prevent an unambiguous behavioural expectation. The abstract
condition \(\Phi_p(x)\) is therefore operationally defined as
\begin{equation}
    \Phi_p(x)
    =
    \underbrace{C_p(x)}_{\text{\scriptsize Applicability}}
    \land
    \underbrace{A_p(x)}_{\text{\scriptsize Admissibility}}
    \land
    \underbrace{\neg E_p(x)}_{\text{\scriptsize Non-exclusion}}.
    \label{eq:operational_conditions}
\end{equation}
The three conditions play complementary roles but need not all be non-trivial
for every property. Depending on the behaviour being specified, admissibility
may follow directly from applicability, or no explicit exclusion condition
may be required.

Let \(P_p\) denote the executable property associated with
behavioural specification \(p\). With \(\Phi_p(x)\) operationalized,
\(P_p\) can be expressed as
\begin{equation}
\begin{aligned}
P_p\bigl(x,D(x)\bigr)
&=
\underbrace{\Phi_p(x)}_{\text{\scriptsize WHEN}}
\Rightarrow
\underbrace{B_p\bigl(x,D(x)\bigr)}_{\text{\scriptsize WHAT}}
\\[1mm]
&=
\underbrace{
\left(
C_p(x)\land A_p(x)\land\neg E_p(x)
\right)
}_{\text{\scriptsize WHEN: operational conditions}}
\Rightarrow
\underbrace{
B_p\bigl(x,D(x)\bigr)
}_{\text{\scriptsize WHAT: expected SUT behaviour}}.
\end{aligned}
\label{eq:executable_property}
\end{equation}
The resulting property explicitly separates the operational conditions under
which \(p\) applies from the expected observable SUT behaviour. Testing \(P_p\) therefore requires generating concrete inputs that satisfy
these operational conditions.

\subsection{Property-Based Test Case Generation}
\label{sec:test_generation}

Test case generation addresses the \emph{when} component of \(P_p\), that is,
the systematic generation of inputs for which the property can be meaningfully
tested. This component is represented by \(\Phi_p(x)\), which must evaluate
to \(\mathsf{true}\) for a generated input to constitute a valid test case.
For this purpose, the framework associates \(P_p\) with a property-specific
generator \(G_p\) defined over a generation space \(\Theta_p\):
\begin{equation}
    G_p : \Theta_p \rightarrow \mathcal{X},
    \qquad
    x = G_p(\theta), \quad \theta \in \Theta_p,
    \label{eq:property_generator}
\end{equation}
where \(\Theta_p\) captures the dimensions along which generated inputs may
vary, including structural, temporal, contextual, or relational variations.

A generated candidate \(x\) constitutes a valid test input for \(P_p\) when
\begin{equation}
    \Phi_p(x)
    =
    \underbrace{C_p(x)}_{\text{\scriptsize Applicability}}
    \land
    \underbrace{A_p(x)}_{\text{\scriptsize Admissibility}}
    \land
    \underbrace{\neg E_p(x)}_{\text{\scriptsize Non-exclusion}}
    =
    \mathsf{true}.
    \label{eq:generation_constraints}
\end{equation}
The constraint may be satisfied by construction,
\(\Phi_p(G_p(\theta))=\mathsf{true}\), or enforced by filtering generated
candidates according to Eq.~\ref{eq:generation_constraints}. Valid test
inputs are submitted to the SUT, whose outputs are evaluated against the
behavioural expectation of \(P_p\) through the property-specific partial
oracle defined next.

\subsection{Property-Specific Partial Oracles}
\label{sec:partial_oracles}

The property-specific partial oracle addresses the \emph{what} component of
\(P_p\), that is, the evaluation of the behavioural expectation represented
by \(B_p(x,D(x))\). Once an input \(x\) is submitted to the SUT and the
observable output \(D(x)\) is obtained, the oracle determines whether this
output satisfies the behavioural expectation of \(P_p\) whenever its
operational conditions hold, rather than whether \(D(x)\) is correct in
general. Let \(O_p\) denote the property-specific partial oracle associated with
\(P_p\):
\begin{equation}
    O_p :
    \mathcal{X} \times \mathcal{Y}
    \rightarrow
    \{\mathsf{pass},\mathsf{fail},\mathsf{undefined}\}.
    \label{eq:partial_oracle}
\end{equation}
Its evaluation is defined as
\begin{equation}
O_p\bigl(x,D(x)\bigr)
=
\begin{cases}
\mathsf{pass},
&
\underbrace{\Phi_p(x)}_{\text{\scriptsize WHEN}}
\land
\underbrace{B_p\bigl(x,D(x)\bigr)}_{\text{\scriptsize WHAT}},
\\[1mm]
\mathsf{fail},
&
\underbrace{\Phi_p(x)}_{\text{\scriptsize WHEN}}
\land
\underbrace{\neg B_p\bigl(x,D(x)\bigr)}_{\text{\scriptsize NOT WHAT}},
\\[1mm]
\mathsf{undefined},
&
\underbrace{\neg \Phi_p(x)}_{\text{\scriptsize NOT WHEN}}.
\end{cases}
\label{eq:partial_oracle_evaluation}
\end{equation}
A \(\mathsf{fail}\) outcome identifies a valid test input for which the
\emph{when} component holds but the \emph{what} component is violated. Such
an input constitutes a counterexample to \(P_p\). In contrast, the
\(\mathsf{undefined}\) outcome reflects the partial nature of \(O_p\): when
the \emph{when} component does not hold, \(\Phi_p(x)=\mathsf{false}\), the
behavioural expectation represented by the \emph{what} component cannot be
meaningfully evaluated for that input.

The complete oracle \(O\) and the property-specific partial oracle \(O_p\)
have fundamentally different semantics. Whereas \(O(x,D(x))\) determines
whether the detector output is correct for input \(x\), \(O_p(x,D(x))\)
determines only whether that output satisfies the behavioural expectation of
\(P_p\) under its operational conditions. Thus, \(O_p\) neither reconstructs
nor approximates \(O\). In particular,
\begin{equation}
    O_p\bigl(x,D(x)\bigr)=\mathsf{pass}
    \;\not\Rightarrow\;
    O\bigl(x,D(x)\bigr)=\mathsf{true}.
    \label{eq:partial_vs_complete_oracle}
\end{equation}
Therefore, a \(\mathsf{pass}\) outcome from \(O_p\) provides evidence of
conformance to a specific executable property, but not of complete output
correctness.

\subsection{Property Checking and Counterexample Discovery}
\label{sec:property_checking}

Property checking combines the \emph{when} and \emph{what} components of
\(P_p\) to evaluate the property over systematically generated valid test
inputs. For each generation configuration \(\theta \in \Theta_p\), the
generator produces an input \(x=G_p(\theta)\). Inputs satisfying
\(\Phi_p(x)=\mathsf{true}\) are submitted to the SUT, and the resulting
output \(D(x)\) is evaluated by the property-specific partial oracle \(O_p\).
Thus, for valid test inputs, a property check can be represented as
\begin{equation}
\begin{gathered}
\underbrace{
    \theta
    \xrightarrow{G_p}
    x
}_{\text{\scriptsize Test generation: WHEN}}
\quad
\underbrace{
    \xrightarrow{D}
    D(x)
}_{\text{\scriptsize SUT execution}}
\quad
\underbrace{
    \xrightarrow{O_p}
    \{\mathsf{pass},\mathsf{fail}\}
}_{\text{\scriptsize Property evaluation: WHAT}}
\end{gathered}
\label{eq:property_check}
\end{equation}
The \(\mathsf{undefined}\) outcome remains part of the general semantics of
\(O_p\) for inputs outside the operational conditions of \(P_p\), as defined
in Eq.~\ref{eq:partial_oracle_evaluation}, but does not arise for valid test
inputs satisfying \(\Phi_p(x)=\mathsf{true}\).

Counterexample discovery follows directly from a failed property check.
The set of counterexamples discovered for \(P_p\) is defined as
\begin{equation}
    \mathcal{CE}_p
    =
    \left\{
        x \in \mathcal{X}
        \mid
        \underbrace{\Phi_p(x)}_{\text{\scriptsize WHEN}}
        \land
        \underbrace{\neg B_p\bigl(x,D(x)\bigr)}_{\text{\scriptsize NOT WHAT}}
    \right\}.
    \label{eq:counterexample_set}
\end{equation}
Equivalently, by the definition of the property-specific partial oracle,
\begin{equation}
    x \in \mathcal{CE}_p
    \quad\Longleftrightarrow\quad
    O_p\bigl(x,D(x)\bigr)=\mathsf{fail}.
    \label{eq:counterexample_oracle}
\end{equation}
Thus, counterexample discovery identifies concrete inputs for which the
operational conditions of \(P_p\) hold but the SUT violates its behavioural
expectation.

The discovered counterexamples provide direct evidence of behavioural
violations under explicitly defined testing conditions. They can therefore be
retained for subsequent analysis, characterization of failure patterns, and
regression testing, without requiring complete instance-level ground truth
over the operational input space.

\begin{figure*}[t]
    \centering
    \includegraphics[width=\textwidth]{figures/PBT_Framework_Final_MaxResolution_600dpi.png}
    \caption{Overview of the proposed property-based testing framework.}
    \label{fig:pbt_framework}
\end{figure*}

Figure~\ref{fig:pbt_framework} summarizes the complete testing process introduced in this section. Domain knowledge is operationalized into executable properties that define
both the conditions under which generated inputs constitute valid test cases
(\emph{when}) and the behavioural expectations against which observable SUT
outputs are evaluated (\emph{what}). Property checking integrates both components through systematic test generation, SUT execution, and property-specific partial-oracle evaluation, with failed evaluations yielding concrete counterexamples.

\section{Empirical Study Design}
\label{sec:empirical_study}

This section presents the empirical instantiation of the PBT framework
introduced in Section~\ref{sec:framework}. The following subsections define
the research questions, describe the dataset and domain context, characterize
the SUT, and present the property suite, PBT configuration, and experimental
protocol.

\subsection{Research Questions}
\label{sec:research_questions}

The empirical study addresses the following research questions:

\begin{itemize}
    \item \textbf{RQ1 -- Property Operationalization:}
    To what extent can domain knowledge about sensor anomalies be
    operationalized as executable properties for testing the ML-based SUT?

 \item \textbf{RQ2 -- Property-Based Behavioural Evaluation:}
How does the ML-based SUT behave across the property-specific generation
space exercised by each executable property?

\item \textbf{RQ3 -- Within-space Behavioural Characterization:}
How does SUT conformance vary across predefined dimensions of the
property-specific generation spaces?

\end{itemize}


\subsection{Empirical, Behavioural, and Operational Context}
\label{sec:dataset_context}

The experimental setting is the AgroConnect\footnote{\url{https://agroconnect.es/}}
greenhouse facility at the University of Almer\'ia (Spain). The facility
covers 1,840~m$^2$ and comprises two independently controlled sectors. Its
instrumentation provides simultaneous measurements of external environmental
conditions, internal greenhouse conditions, and climate-control actions.
Figure~\ref{fig:agroconnect_setup} shows the experimental greenhouse and
sensing environment.

\begin{figure}[h!]
\centering
\begin{subfigure}{0.39\linewidth}
    \centering
    \includegraphics[width=\linewidth]
    {figures/agroconnect_greenhouse_large.jpeg}
    \caption{Experimental greenhouse facility.}
    \label{fig:agro_overview}
\end{subfigure}
\hfill
\begin{subfigure}{0.28\linewidth}
    \centering
    \includegraphics[width=\linewidth]{figures/agro_crop.jpeg}
    \caption{Monitored crop environment.}
    \label{fig:agro_crop}
\end{subfigure}
\hfill
\begin{subfigure}{0.21\linewidth}
    \centering
    \includegraphics[width=\linewidth]{figures/agroconnect_station.png}
    \caption{IoT sensing infrastructure.}
    \label{fig:agro_station}
\end{subfigure}
\caption{AgroConnect experimental greenhouse and sensing environment.}
\label{fig:agroconnect_setup}
\end{figure}

The empirical basis of the study is a temporally ordered multivariate dataset
collected from the AgroConnect infrastructure. The dataset spans from
August 15, 2019 to May 31, 2022, with measurements recorded at a fixed
sampling interval of 30~s, resulting in 2,940,480 observations. The monitored
variables represent three complementary dimensions of greenhouse operation:
external environmental conditions, internal greenhouse state, and control
actions. External conditions include CO$_2$ concentration, relative humidity,
global solar radiation, air temperature, and wind speed. The internal
greenhouse state comprises CO$_2$ concentration, relative humidity, global
solar radiation, and air temperature, while roof and side ventilation
openings characterize the operational state of the greenhouse control system.
The prefixes \textbf{D}, \textbf{X}, and \textbf{U} denote disturbance
variables, state variables, and control signals, respectively.
Table~\ref{tab:descriptive_stats} summarizes the monitored variables and
their descriptive statistics.

\begin{table*}[t]
\caption{Descriptive statistics of the greenhouse variables used in the
empirical study.}
\label{tab:descriptive_stats}
\centering
{\scriptsize
\resizebox{\textwidth}{!}{%
\begin{tabular}{|c|c|c|c|c|c|c|}
\hline
\textbf{Variable} &
\textbf{Description} &
\textbf{Mean} &
\textbf{Std. Dev.} &
\textbf{Min} &
\textbf{Max} &
\textbf{Median} \\
\hline
\hline
{\tt DCO2EXT} &
Outdoor CO$_2$ concentration (ppm) &
428.28 & 12.34 & 339.92 & 522.23 & 428.79 \\

{\tt DHEXT} &
Outdoor relative humidity (\%) &
63.31 & 16.75 & 10.27 & 98.06 & 64.03 \\

{\tt DGREXT} &
Outdoor global solar radiation (W/m$^2$) &
182.99 & 272.20 & 0.00 & 1473.48 & 7.37 \\

{\tt DTEXT} &
Outdoor air temperature ($^\circ$C) &
17.41 & 5.62 & 3.24 & 35.40 & 16.57 \\

{\tt DVEXT} &
Outdoor wind speed (m/s) &
2.24 & 2.20 & 0.00 & 23.59 & 1.47 \\
\hline

{\tt XCO2INT} &
Indoor CO$_2$ concentration (ppm) &
406.70 & 42.50 & 210.94 & 721.04 & 401.50 \\

{\tt XHINT} &
Indoor relative humidity (\%) &
69.80 & 18.99 & 11.57 & 96.85 & 76.60 \\

{\tt XGRINT} &
Indoor global solar radiation (W/m$^2$) &
83.14 & 131.43 & 0.00 & 990.22 & 1.36 \\

{\tt XTINT} &
Indoor air temperature ($^\circ$C) &
19.69 & 7.15 & 5.97 & 48.19 & 18.57 \\
\hline

{\tt UVENT\_CEN} &
Roof vent opening (\%) &
34.71 & 44.34 & 0.00 & 100.00 & 0.00 \\

{\tt UVENT\_INT} &
Side vent opening (\%) &
37.76 & 45.17 & 0.00 & 100.00 & 0.00 \\
\hline
\end{tabular}
}
}

\vspace{-3mm}
\begin{center}
\footnotesize
\textit{Note:} In variable names, the prefix {\bf D} denotes disturbance
variables, {\bf X} state variables, and {\bf U} control signals.
\end{center}
\end{table*}

The behavioural context is determined by the temporal, multivariate, and
context-dependent nature of greenhouse climate dynamics. Individual
measurements cannot always be interpreted in isolation, because their
meaning may depend on recent temporal evolution, external environmental
conditions, control actions, and relationships among monitored variables.
These temporal and multivariate dependencies provide the behavioural basis
for the SUT representation formalized in Section~\ref{sec:sut}.

This empirical and behavioural characterization also underlies the anomaly
taxonomy established in our previous work
\citep{carrero2026ifac,carrero2026qrs}. Recurrent abnormal behaviours
identified through exploratory analysis and domain interpretation were
organized into the anomaly categories subsequently used by the operational
pipeline. This taxonomy provides the domain-grounded semantic knowledge
formalized in Section~\ref{sec:behavioural_knowledge}.

The operational context is provided by the ML-based anomaly detection and
semantic characterization pipeline developed in our previous work. As shown
in Figure~\ref{fig:anomaly_pipeline}, temporally ordered measurements are
transformed into a feature representation and processed through anomaly
detection and semantic characterization stages. The pipeline combines two
ML components implemented using LightGBM with deterministic conditions for
directly identifiable situations, including missing measurements,
out-of-range values, and system-level failures. Its concrete representation
as the SUT is defined in the next subsection.

\begin{figure*}[t]
    \centering
    \includegraphics[width=\textwidth]
    {figures/operational_anomaly_pipeline_PBT_taxonomy_linked.pdf}
    \caption{Operational anomaly detection and semantic characterization
    pipeline considered in the empirical study.}
    \label{fig:anomaly_pipeline}
\end{figure*}


\subsection{System Under Test}
\label{sec:sut}

The empirical SUT instantiates the abstract mapping
\(D:\mathcal{X}\rightarrow\mathcal{Y}\) introduced in
Section~\ref{sec:problem_formulation}. Its input space \(\mathcal{X}\)
comprises temporally ordered multivariate greenhouse observations, while
\(D\) encapsulates the complete anomaly-processing pipeline described in
Section~\ref{sec:dataset_context}.

An input \(x\in\mathcal{X}\) is represented as a temporally ordered
multivariate sequence of greenhouse observations,
\begin{equation}
    x =
    \bigl(
        (t_1,\mathbf{v}_1),
        (t_2,\mathbf{v}_2),
        \ldots,
        (t_n,\mathbf{v}_n)
    \bigr),
    \label{eq:sut_input}
\end{equation}
where \(t_i\) denotes the timestamp of the \(i\)-th observation and
\(\mathbf{v}_i\) is the multivariate observation vector at that timestamp,
defined as
\begin{equation}
    \mathbf{v}_i =
    \bigl(
        v_{i,1}, v_{i,2}, \ldots, v_{i,m}
    \bigr),
    \label{eq:sut_observation}
\end{equation}
where \(v_{i,j}\) denotes the observed value of the \(j\)-th monitored
variable at timestamp \(t_i\), and \(m\) is the number of variables
represented in the operational input.

The input representation supports two complementary views of a target
variable. For example, if indoor air temperature (\texttt{XTINT}) corresponds
to the \(k\)-th component of \(\mathbf{v}_i\), its value at timestamp \(t_i\)
can be represented within the multivariate observation as
$$
    \mathbf{v}_i =
    \bigl(
        v_{i,1},\ldots,
        \underbrace{v_{i,k}}_{\texttt{XTINT}_i},
        \ldots,v_{i,m}
    \bigr),
$$
providing a multivariate view in which the target value is considered
together with the simultaneous values of the remaining monitored variables.
Conversely, following the same component across consecutive timestamps,
$$
    \bigl(
        \underbrace{v_{1,k}}_{\texttt{XTINT}_1},
        \underbrace{v_{2,k}}_{\texttt{XTINT}_2},
        \ldots,
        \underbrace{v_{n,k}}_{\texttt{XTINT}_n}
    \bigr),
$$
provides a temporal view of its evolution. These complementary views capture
the cross-variable relationships and temporal behaviour required by the
anomaly-processing pipeline.

The mapping \(D\) encapsulates the complete anomaly-processing pipeline shown
in Figure~\ref{fig:anomaly_pipeline}, including feature construction and
preprocessing, anomaly detection, decision logic, and semantic
characterization. These operations remain internal to the SUT; in particular,
the feature representation constructed from \(x\) is part of the execution of
\(D\), rather than the externally specified input space \(\mathcal{X}\).

The observable SUT response is defined as
\begin{equation}
    D(x) =
    \bigl(
        d(x), c(x)
    \bigr),
    \label{eq:sut_output}
\end{equation}
where \(d(x)\) denotes the final anomaly decision and \(c(x)\) the
corresponding semantic category when an anomaly is detected. This
input--output boundary provides the interface through which the executable
properties evaluate the externally observable behaviour of the SUT.


\subsection{Domain-Grounded Behavioural Knowledge}
\label{sec:behavioural_knowledge}
The anomaly taxonomy established in our previous work provides the
domain-grounded behavioural knowledge used to instantiate the executable
properties in this study. Different anomaly behaviours require different
forms of evidence within the SUT input \(x\), depending on whether they can
be characterized from an individual observation, its temporal evolution, or
its relationships with other monitored variables.

To distinguish these forms of evidence, the \emph{evidence scope} identifies
the part of \(x\) that must be considered to characterize each behaviour.
Table~\ref{tab:evidence_scope} defines the three evidence scopes used in this
study, together with their representation within \(x\) and their semantic
interpretation.

\begin{table}[t]
\caption{Semantic definition of the evidence scopes used to organize
behavioural knowledge.}
\label{tab:evidence_scope}
\centering
\scriptsize
\begin{tabular}{
>{\centering\arraybackslash}p{0.10\linewidth}
p{0.32\linewidth}
p{0.48\linewidth}
}
\hline
\textbf{Scope} &
\textbf{Representation} &
\textbf{Semantic interpretation} \\
\hline

\(\mathrm{O}\) &
\(v_{i,k}\) &
Observation-level evidence for which an individual value of the target
variable \(k\) at timestamp \(t_i\) is sufficient for behavioural
characterization. \\

\(\mathrm{T}\) &
\(\displaystyle
(v_{i,k},v_{i+1,k},\ldots,v_{i+r,k})
\) &
Temporal evidence for which behavioural characterization requires the
evolution of the target variable \(k\) over a local window of \(r+1\)
consecutive observations, extracted from the \emph{vertical view}
introduced in Section~\ref{sec:sut}. \\

\(\mathrm{M}\) &
\(\displaystyle
\bigl(
v_{i,1},\ldots,
\underbrace{v_{i,k}}_{\mathrm{target}},
\ldots,v_{i,m}
\bigr)
\) &
Multivariate evidence for which behavioural characterization requires
the simultaneous relationships between the target variable and the
remaining monitored variables at timestamp \(t_i\), corresponding to
the \emph{horizontal view} introduced in Section~\ref{sec:sut}. \\
\hline
\end{tabular}
\end{table}
Based on this distinction, the domain knowledge associated with each
behaviour can be described systematically in terms of its behavioural
characterization, the evidence required to support that characterization,
and the domain constraints governing its interpretation.
Table~\ref{tab:knowledge_map} provides this knowledge map for the anomaly
categories considered in the study.

\begin{table}[!t]
\caption{Domain-grounded knowledge map for the anomaly behaviours.}
\label{tab:knowledge_map}
\centering
\scriptsize
\setlength{\tabcolsep}{2pt}
\renewcommand{\arraystretch}{1.12}

\begin{tabular}{
>{\centering\arraybackslash}m{0.14\linewidth}
p{0.38\linewidth}
>{\centering\arraybackslash}m{0.08\linewidth}
p{0.32\linewidth}
}
\hline

\multicolumn{1}{c}{\textbf{Behaviour}} &
\multicolumn{1}{c}{\textbf{Behavioural characterization}} &
\multicolumn{1}{c}{\shortstack{\textbf{Evidence}\\\textbf{scope}}} &
\multicolumn{1}{c}{\textbf{Domain constraints}} \\

\hline

\textbf{MV (Missing Values)} &
Explicit absence of a monitored value, represented as \texttt{NaN}. &
\(\mathrm{O}\) &
Measurements must be available for every monitored variable at each
timestamp of the regularly sampled sequence. \\[1mm]

\textbf{SS (Stuck Sensor)} &
Persistence of the same or minimally varying value over consecutive
timestamps, producing an abnormally constant sequence for the monitored
variable. &
\(\mathrm{T}\) &
Temporal variation must remain consistent with the variable-specific
dynamics; prolonged constancy beyond the admissible duration is excluded
under normal operation. \\[1mm]

\textbf{N (Noise)} &
Erratic fluctuations, spikes, or sudden drops that depart from the local
temporal dynamics of the monitored variable. &
\(\mathrm{T}\) &
Short-term changes must remain within the variable-specific temporal
dynamics and admissible rates of change. \\[1mm]

\textbf{ORV (Out-of-Range Values)} &
Observed values lying outside the admissible range established for the
monitored variable. &
\(\mathrm{O}\) &
Values must remain within the variable-specific physical or operational
bounds. \\[1mm]

\textbf{CA (Contextual Anomalies)} &
Observations that do not conform to an admissible contextual configuration
of the monitored greenhouse system, despite being individually plausible. &
\(\mathrm{M}\) &
Values and variable configurations must remain compatible with the
applicable greenhouse operating context, such as diurnal or nocturnal
conditions and open or closed ventilation states. \\[1mm]

\textbf{CD (Correlation Deviations)} &
Deviations of the observed relationship between monitored variables from
their expected correlation structure. &
\(\mathrm{M}\) &
Inter-variable relationships must remain consistent with the expected
correlation structure. \\[1mm]

\textbf{PSF (Partial System Failure)} &
Simultaneous loss or invalid acquisition of multiple monitored values while
the monitored system remains partially available. &
\(\mathrm{M}\) &
Loss or invalid acquisition may affect only a subset of monitored variables;
complete system unavailability is excluded. \\

\hline
\end{tabular}
\end{table}

Figure~\ref{fig:evidence_scope_examples} illustrates the three evidence
scopes through representative anomaly behaviours observed in real greenhouse
data: ORV in XGRINT as observation-level evidence (\(\mathrm{O}\)), SS in
XTINT as temporal evidence (\(\mathrm{T}\)), and CA in DGREXT under nighttime
operating conditions as multivariate evidence (\(\mathrm{M}\)).

\begin{figure*}[!t]
    \centering
    % Observation-level evidence (O)
    \begin{subfigure}[t]{0.48\textwidth}
        \centering
        \includegraphics[width=\linewidth]
        {figures/04b_out_of_range_XGRINT_multiple_validation.png}
        \caption{Out-of-range values (ORV) in XGRINT:
        observation-level evidence (\(\mathrm{O}\)).}
        \label{fig:example_orv}
    \end{subfigure}
    \hfill
    % Temporal evidence (T)
    \begin{subfigure}[t]{0.48\textwidth}
        \centering
        \includegraphics[width=\linewidth]
        {figures/02_sensor_atascado_XTINV_validacion.png}
        \caption{Stuck-sensor behaviour (SS) in XTINT:
        temporal evidence (\(\mathrm{T}\)).}
        \label{fig:example_ss}
    \end{subfigure}

    \vspace{0.5em}

    % Multivariate/contextual evidence (M)
    \begin{subfigure}[t]{0.62\textwidth}
        \centering
        \includegraphics[width=\linewidth]
        {figures/05b_contextual_DGREXT_night_radiation_real.png}
      \caption{Contextual anomaly (CA) in DGREXT under nighttime
operating conditions: multivariate evidence (\(\mathrm{M}\)).}
        \label{fig:example_ca}
    \end{subfigure}
        \caption{Representative anomaly behaviours for the three evidence scopes.}
    \label{fig:evidence_scope_examples}
\end{figure*}

The resulting knowledge map provides the basis for the property suite
introduced in the following subsection, where these behavioural
characterizations and domain constraints are operationalized as executable
properties, property-specific generators, and partial behavioural oracles.


\subsection{Property Suite}
\label{sec:property_suite}

The domain-grounded knowledge established in the previous subsection is
instantiated as seven executable properties following the PBT framework
introduced in Section~\ref{sec:framework}. For each property \(P_p\), the
applicability, admissibility, and exclusion conditions define
\[
\Phi_p(x)=C_p(x)\land A_p(x)\land\neg E_p(x),
\]
which specifies \emph{when} the property applies, while \(B_p(x,D(x))\)
specifies \emph{what} behaviour is expected from the observable SUT response
\(D(x)=(d(x),c(x))\). The properties are derived from the prior empirical
anomaly characterization and domain constraints; SUT predictions and PBT
counterexamples are not used in their definition. The seven executable properties are
instantiated as follows.
\begingroup
\scriptsize
\begin{equation}
\begin{aligned}
P_{\mathrm{MV}}\bigl(x,D(x)\bigr)
&=
\overbrace{
    \underbrace{v_{i,k}}_{\mathrm{O}}=\texttt{NaN}
}^{C_{\mathrm{MV}}(x)}
\;\land\;
\overbrace{
    \top
}^{A_{\mathrm{MV}}(x)}
\;\land\;
\overbrace{
    \text{no PSF condition holds}
}^{\neg E_{\mathrm{MV}}(x)}
\;\Rightarrow\;
\overbrace{
    d(x)=\mathrm{anomaly}
    \;\land\;
    c(x)=\mathrm{MV}
}^{B_{\mathrm{MV}}(x,D(x))}.
\end{aligned}
\label{eq:property_mv}
\end{equation}
\endgroup

The non-exclusion condition prevents an individual missing value from being
interpreted as MV when it belongs to a simultaneous missingness pattern
characterized as PSF.
\begingroup
\scriptsize
\begin{equation}
\begin{aligned}
P_{\mathrm{SS}}\bigl(x,D(x)\bigr)
&=
\overbrace{
    \operatorname{Var}\!\left(
        \underbrace{
        (v_{i,k},v_{i+1,k},\ldots,v_{i+r,k})
        }_{\mathrm{T}}
    \right)
    \approx 0
}^{C_{\mathrm{SS}}(x)}
\;\land\;
\overbrace{
    \top
}^{A_{\mathrm{SS}}(x)}
\;\land\;
\overbrace{
    \text{no MV, ORV, or PSF condition holds}
}^{\neg E_{\mathrm{SS}}(x)}
\\[1mm]
&\Rightarrow
\overbrace{
    d(x)=\mathrm{anomaly}
    \;\land\;
    c(x)=\mathrm{SS}
}^{B_{\mathrm{SS}}(x,D(x))}.
\end{aligned}
\label{eq:property_ss}
\end{equation}
\endgroup
The non-exclusion condition prevents a near-zero-variance temporal pattern
from being interpreted as SS when an MV, ORV, or PSF condition provides an
alternative anomaly interpretation.
\begingroup
\scriptsize
\begin{equation}
\begin{aligned}
P_{\mathrm{N}}\bigl(x,D(x)\bigr)
&=
\overbrace{
    \substack{
    \left|v_{i,k}-v_{i-1,k}\right|>\delta_k
    \;\land\;
    \left|v_{i,k}-v_{i+1,k}\right|>\delta_k\\[1mm]
       \underbrace{
    (v_{i-1,k},\,v_{i,k},\,v_{i+1,k})
    }_{\mathrm{T}}
    }
}^{C_{\mathrm{N}}(x)}
\;\land\;
\overbrace{
    L_k\leq v_{i,k}\leq U_k
}^{A_{\mathrm{N}}(x)}
\;\land\;
\overbrace{
    \text{no PSF condition holds}
}^{\neg E_{\mathrm{N}}(x)}
\\[1mm]
&\Rightarrow
\overbrace{
    d(x)=\mathrm{anomaly}
    \;\land\;
    c(x)=\mathrm{N}
}^{B_{\mathrm{N}}(x,D(x))}.
\end{aligned}
\label{eq:property_n}
\end{equation}
\endgroup
\(\delta_k\) denotes the sensor-specific deviation threshold, while
\(L_k\) and \(U_k\) denote the corresponding lower and upper physical
limits. The admissibility condition keeps the target value within its valid
physical range, while the non-exclusion condition rules out concurrent PSF
patterns that would prevent an unambiguous N interpretation.
\begingroup
\scriptsize
\begin{equation}
\begin{aligned}
P_{\mathrm{ORV}}\bigl(x,D(x)\bigr)
&=
\overbrace{
    \underbrace{v_{i,k}}_{\mathrm{O}}\notin[L_k,U_k]
}^{C_{\mathrm{ORV}}(x)}
\;\land\;
\overbrace{
    \top
}^{A_{\mathrm{ORV}}(x)}
\;\land\;
\overbrace{
    \text{no PSF condition holds}
}^{\neg E_{\mathrm{ORV}}(x)}
\\[1mm]
&\Rightarrow
\overbrace{
    d(x)=\mathrm{anomaly}
    \;\land\;
    c(x)=\mathrm{ORV}
}^{B_{\mathrm{ORV}}(x,D(x))}.
\end{aligned}
\label{eq:property_orv}
\end{equation}
\endgroup
\(L_k\) and \(U_k\) denote the lower and upper physical limits of sensor
\(k\). The non-exclusion condition rules out concurrent PSF patterns that
would prevent an unambiguous ORV interpretation.

\begingroup
\scriptsize
\begin{equation}
\begin{aligned}
P_{\mathrm{CA}}\bigl(x,D(x)\bigr)
&=
\overbrace{
    \underbrace{
    (v_{i,1},\ldots,v_{i,k},\ldots,v_{i,m})
    }_{\mathrm{M}}
    \text{ exhibits contextual incompatibility}
}^{C_{\mathrm{CA}}(x)}
\\[1mm]
&\land\;
\overbrace{
 \forall k\in\{1,\ldots,m\}:
L_k\leq v_{i,k}\leq U_k
}^{A_{\mathrm{CA}}(x)}
\;\land\;
\overbrace{
    \text{no PSF condition holds}
}^{\neg E_{\mathrm{CA}}(x)}
\\[1mm]
&\Rightarrow
\overbrace{
    d(x)=\mathrm{anomaly}
    \;\land\;
    c(x)=\mathrm{CA}
}^{B_{\mathrm{CA}}(x,D(x))}.
\end{aligned}
\label{eq:property_ca}
\end{equation}
\endgroup
Contextual incompatibility refers to violations of the relationships
governing radiation consistency, indoor--outdoor temperature conditions,
or indoor CO$_2$ under operating contexts representing diurnal or nocturnal
conditions and open or closed ventilation states. The admissibility condition
ensures that all monitored values remain individually within their physical
limits, while the non-exclusion condition rules out PSF patterns that would
prevent an unambiguous CA interpretation.

\begingroup
\scriptsize
\begin{equation}
\begin{aligned}
P_{\mathrm{CD}}\bigl(x,D(x)\bigr)
&=
\overbrace{
    \substack{
    \exists\,\ell\neq k:
    Q_{k\ell}(v_{i,k},v_{i,\ell})
    \\[1mm]
    \underbrace{
    (v_{i,1},\ldots,v_{i,k},\ldots,v_{i,\ell},\ldots,v_{i,m})
    }_{\mathrm{M}}
    }
}^{C_{\mathrm{CD}}(x)}
\\[1mm]
&\land\;
\overbrace{
\forall k\in\{1,\ldots,m\}:
L_k\leq v_{i,k}\leq U_k
}^{A_{\mathrm{CD}}(x)}
\;\land\;
\overbrace{
    \text{no PSF condition holds}
}^{\neg E_{\mathrm{CD}}(x)}
\\[1mm]
&\Rightarrow
\overbrace{
    d(x)=\mathrm{anomaly}
    \;\land\;
    c(x)=\mathrm{CD}
}^{B_{\mathrm{CD}}(x,D(x))}.
\end{aligned}
\label{eq:property_cd}
\end{equation}
\endgroup

where

\begingroup
\scriptsize
\begin{equation}
Q_{k\ell}(v_{i,k},v_{i,\ell})=
\begin{cases}
(v_{i,\ell}>p_{75,\ell}\land v_{i,k}\leq p_{25,k})
\;\lor\;
(v_{i,\ell}<p_{25,\ell}\land v_{i,k}\geq p_{75,k}),
& \rho_{k\ell}>0,\\[1mm]
(v_{i,\ell}>p_{75,\ell}\land v_{i,k}\geq p_{75,k})
\;\lor\;
(v_{i,\ell}<p_{25,\ell}\land v_{i,k}\leq p_{25,k}),
& \rho_{k\ell}<0.
\end{cases}
\label{eq:cd_incompatibility}
\end{equation}
\endgroup

\(\rho_{k\ell}\) denotes the empirical correlation between variables
\(k\) and \(\ell\), while \(p_{25,j}\) and \(p_{75,j}\) denote the
empirical 25th and 75th percentiles of variable \(j\).
The predicate \(Q_{k\ell}\) captures pairwise incompatibility with the
expected correlation structure: for positively correlated variables,
opposite extreme quartiles characterize the deviation, whereas for
negatively correlated variables, same-side extreme quartiles characterize
it. The non-exclusion condition rules out PSF patterns that would prevent
an unambiguous CD interpretation.

\begingroup
\scriptsize
\begin{equation}
\begin{aligned}
P_{\mathrm{PSF}}\bigl(x,D(x)\bigr)
&=
\overbrace{
    \substack{
    \left|
    \{k:v_{i,k}=\texttt{NaN}\}
    \right|\geq\eta
    \\[1mm]
    \underbrace{
    (v_{i,1},\ldots,v_{i,k},\ldots,v_{i,m})
    }_{\mathrm{M}}
    }
}^{C_{\mathrm{PSF}}(x)}
\;\land\;
\overbrace{
    \top
}^{A_{\mathrm{PSF}}(x)}
\;\land\;
\overbrace{
    \exists\,k\in\{1,\ldots,m\}:
    v_{i,k}\neq\texttt{NaN}
}^{\neg E_{\mathrm{PSF}}(x)}
\\[1mm]
&\Rightarrow
\overbrace{
    d(x)=\mathrm{anomaly}
    \;\land\;
    c(x)=\mathrm{PSF}
}^{B_{\mathrm{PSF}}(x,D(x))}.
\end{aligned}
\label{eq:property_psf}
\end{equation}
\endgroup
\(\eta\) denotes the minimum number of simultaneously missing variables
required to characterize a partial system failure. The non-exclusion
condition requires at least one monitored variable to remain available,
thereby distinguishing PSF from complete system failure.


\begin{table}[!t]
\caption{Experimental configuration of the property-specific generation
spaces \(\Theta_p\).}
\label{tab:pbt_generation_config}
\centering

{\scriptsize
\setlength{\tabcolsep}{3pt}
\begin{tabular}{|c|p{10.7cm}|}
\hline
\textbf{Property \(P_p\)} &
\textbf{Experimental configuration of \(\Theta_p\)} \\
\hline
\hline

\(P_{\mathrm{MV}}\) &
Observation: any \(i\);
variable \(k\): \texttt{DCO2EXT}, \texttt{DHEXT}, \texttt{DGREXT},
\texttt{DTEXT}, \texttt{DVEXT}, \texttt{XCO2INT}, \texttt{XGRINT},
\texttt{XHINT}, or \texttt{XTINT}. \\
\hline

\(P_{\mathrm{SS}}\) &
Temporal window: \((v_{i,k},\ldots,v_{i+r,k})\);
variable \(k\): \texttt{DCO2EXT}, \texttt{DHEXT}, \texttt{DTEXT},
\texttt{DVEXT}, \texttt{XCO2INT}, \texttt{XGRINT}, \texttt{XHINT},
or \texttt{XTINT};
SS-window length (samples):
\texttt{DCO2EXT}: 20--40;
\texttt{DHEXT}: 16--30;
\texttt{DTEXT}: 40--60;
\texttt{DVEXT}: 20--60;
\texttt{XCO2INT}: 20--60;
\texttt{XGRINT}: 20--60;
\texttt{XHINT}: 16--30;
\texttt{XTINT}: 40--60;
\texttt{XGRINT} context:
\texttt{XGRINT} \(>10\),
\texttt{DGREXT} \(>10\), and
\(\operatorname{Std}_{30\mathrm{min}}(\texttt{DGREXT})>2\).\\
\hline

\(P_{\mathrm{N}}\) &
Temporal window: \((v_{i-1,k},v_{i,k},v_{i+1,k})\);
variable \(k\): \texttt{DCO2EXT}, \texttt{DHEXT}, \texttt{DGREXT},
\texttt{DTEXT}, \texttt{DVEXT}, \texttt{XCO2INT}, \texttt{XGRINT},
\texttt{XHINT}, or \texttt{XTINT};
perturbation magnitude:
\texttt{DCO2EXT}: 25--75 ppm;
\texttt{DHEXT}: 10--25\%;
\texttt{DGREXT}: 150--500 W/m$^2$;
\texttt{DTEXT}: 3--8 $^\circ$C;
\texttt{DVEXT}: 5--9 m/s;
\texttt{XCO2INT}: 60--160 ppm;
\texttt{XGRINT}: 80--250 W/m$^2$;
\texttt{XHINT}: 10--25\%;
\texttt{XTINT}: 3--8 $^\circ$C;
direction: positive or negative, except
\texttt{DGREXT}, \texttt{XGRINT}, and \texttt{DVEXT} (positive only). \\
\hline

\(P_{\mathrm{ORV}}\) &
Observation: any \(i\);
variable \(k\): \texttt{DCO2EXT}, \texttt{DHEXT}, \texttt{DGREXT},
\texttt{DTEXT}, \texttt{DVEXT}, \texttt{XCO2INT}, \texttt{XGRINT},
\texttt{XHINT}, or \texttt{XTINT};
physical interval \([L_k,U_k]\):
\texttt{DCO2EXT}: [300,800] ppm;
\texttt{DHEXT}: [0,100]\%;
\texttt{DGREXT}: [0,1500] W/m$^2$;
\texttt{DTEXT}: [-20,60] $^\circ$C;
\texttt{DVEXT}: [0,50] m/s;
\texttt{XCO2INT}: [200,3000] ppm;
\texttt{XGRINT}: [0,1200] W/m$^2$;
\texttt{XHINT}: [0,100]\%;
\texttt{XTINT}: [-5,60] $^\circ$C;
direction: below \(L_k\) or above \(U_k\);
out-of-range margin: \(m_k=\alpha(U_k-L_k)\), with
\(\alpha\in\{0.01,0.025,0.05,0.10\}\). \\
\hline

\hline
\end{tabular}

\vspace{1mm}
\parbox{\textwidth}{
\scriptsize
\textbf{\textit{Stuck note:}} The reported SS-window ranges define the temporal-window
lengths explored during test generation and do not represent lower or upper
semantic thresholds for stuck-sensor behaviour. The \texttt{XGRINT} range
derives from the legacy empirical configuration and is therefore applied
only under an active-radiation context. Specifically, \texttt{XGRINT}>10
and \texttt{DGREXT}>10 ensure that both internal and external radiation are
active, while
\(\operatorname{Std}_{30\mathrm{min}}(\texttt{DGREXT})>2\)
requires external radiation to exhibit sufficient variability over a
30-min sliding window (60 observations at the 30-s sampling interval).
These conditions prevent constant or near-zero radiation regimes, such as
nighttime periods, from being used to generate ambiguous SS test cases.  \texttt{DGREXT} is not included as an SS target variable because persistent constant or near-constant radiation may correspond to physically valid behaviour, and the available empirical evidence does not provide an unambiguous generation configuration for SS testing.\\
\textbf{\textit{Noise note:}} The reported perturbation ranges define the
sensor-specific absolute magnitudes explored when generating N test cases.
The corresponding deviation threshold \(\delta_k\) is instantiated as the
lower bound of each sensor-specific perturbation range:
25 ppm (\texttt{DCO2EXT}), 10\% (\texttt{DHEXT}),
150 W/m$^2$ (\texttt{DGREXT}), 3 $^\circ$C (\texttt{DTEXT}),
5 m/s (\texttt{DVEXT}), 60 ppm (\texttt{XCO2INT}),
80 W/m$^2$ (\texttt{XGRINT}), 10\% (\texttt{XHINT}), and
3 $^\circ$C (\texttt{XTINT}).
Perturbations are generated in both directions except for
\texttt{DGREXT}, \texttt{XGRINT}, and \texttt{DVEXT}, for which only
positive perturbations are considered. Negative radiation deviations may
represent physically valid short-term irradiance reductions, such as those
caused by cloud passage, while negative wind-speed deviations may correspond
to plausible calm conditions.\\
\textbf{\textit{ORV note:}} The out-of-range margin \(m_k\) is a
stratified sensor-specific displacement from the corresponding physical
boundary. Thus, generated values are \(L_k-m_k\) or \(U_k+m_k\), with both
directions and all four \(\alpha\) levels considered for every monitored
variable.\\

}
}

\end{table}


\begin{table}[t]
\ContinuedFloat
\caption{Experimental configuration of the property-specific generation
spaces \(\Theta_p\) (continued).}
\label{tab:pbt_generation_config}
\centering

{\scriptsize
\setlength{\tabcolsep}{3pt}
\begin{tabular}{|c|p{10.7cm}|}
\hline
\textbf{Property \(P_p\)} &
\textbf{Experimental configuration of \(\Theta_p\)} \\
\hline
\hline

\(P_{\mathrm{CA}}\) &
Multivariate observation:
\((v_{i,1},\ldots,v_{i,k},\ldots,v_{i,m})\);
variable \(k\):
\texttt{XGRINT}, \texttt{XCO2INT}, or \texttt{XTINT};
context:
\texttt{XGRINT}/\texttt{DGREXT} radiation consistency,
\texttt{XCO2INT}/ventilation/\texttt{DGREXT} operating conditions,
or \texttt{XTINT}/\texttt{DTEXT}/\texttt{DGREXT} temperature conditions;
generation parameters:
\texttt{XGRINT}: anomalous ratio 0.02--0.08;
\texttt{XCO2INT}: 700--1200 ppm, segment length 3--8 observations,
within-segment variation \(\pm 10\) ppm;
\texttt{XTINT}: temperature difference 3--10 \(^{\circ}\)C. \\
\hline

\(P_{\mathrm{CD}}\) &
Multivariate observation:
\((v_{i,1},\ldots,v_{i,k},\ldots,v_{i,\ell},\ldots,v_{i,m})\);
variable pair \((k,\ell)\):
(\texttt{DGREXT},\texttt{XGRINT}),
(\texttt{XTINT},\texttt{XHINT}),
(\texttt{DTEXT},\texttt{XTINT}), or
(\texttt{DHEXT},\texttt{XHINT});
expected correlation \(\rho_{k\ell}\):
0.957, -0.819, 0.932, and 0.737, respectively;
relationship-selection threshold:
\(|\rho_{k\ell}|>0.3\);
pairwise incompatibility:
reference variable \(v_{i,\ell}\) below \(p_{25,\ell}\) or above
\(p_{75,\ell}\), with target variable \(v_{i,k}\) placed at
\(p_{25,k}\) or \(p_{75,k}\) according to the correlation sign. \\
\hline

\(P_{\mathrm{PSF}}\) &
Multivariate observation:
\((v_{i,1},\ldots,v_{i,k},\ldots,v_{i,m})\);
variable \(k\): \texttt{DCO2EXT}, \texttt{DHEXT}, \texttt{DGREXT},
\texttt{DTEXT}, \texttt{DVEXT}, \texttt{XCO2INT}, \texttt{XGRINT},
\texttt{XHINT}, or \texttt{XTINT};
simultaneously missing variables: 4--8;
minimum missing-variable threshold: \(\eta=4\);
missing value: \texttt{NaN}. \\
\hline

\end{tabular}

\vspace{1mm}
\parbox{\textwidth}{
\scriptsize
\textbf{\textit{CA note:}} The radiation-consistency and temperature-context
configurations are generated at a single multivariate observation, whereas
the CO$_2$ operating-context configuration spans 3--8 consecutive
observations; each observation in the resulting segment is evaluated
against \(P_{\mathrm{CA}}\).\\
\textbf{\textit{CD note:}} \(p_{25,k}\) and \(p_{75,k}\) denote the empirical
25th and 75th percentiles of the target variable, while \(p_{25,\ell}\)
and \(p_{75,\ell}\) denote those of the reference variable. For positively
correlated pairs, a reference value above \(p_{75,\ell}\) is combined
with a target value at \(p_{25,k}\), and a reference value below
\(p_{25,\ell}\) with a target value at \(p_{75,k}\). For negatively
correlated pairs, the target is placed at the same-side quartile.
The threshold \(|\rho_{k\ell}|>0.3\) is used only to select relationships
with sufficient empirical correlation for CD test generation.\\
\textbf{\textit{PSF note:}} The threshold \(\eta=4\) is inherited from the
deterministic PSF rule applied to real missingness patterns in the
preceding anomaly-analysis pipeline. For PBT generation, PSF cases are
instantiated by simultaneously setting 4--8 of the nine monitored
variables to \texttt{NaN}. The upper bound excludes the all-missing
configuration, which corresponds to complete system failure rather than
PSF.
}
}

\end{table}

\subsection{PBT Generation Configuration and Experimental Protocol}
\label{sec:pbt_protocol}

The executable properties defined in Section~\ref{sec:property_suite} were
instantiated through property-specific generation spaces \(\Theta_p\) and
generators \(G_p\), following the test-generation mechanism introduced in
Section~\ref{sec:test_generation}. Table~\ref{tab:pbt_generation_config}
summarizes the experimental configuration adopted for each property.

The generation spaces were instantiated from the empirical characterization,
domain knowledge, and sensor-specific parameterizations established in the
preceding anomaly-analysis pipeline~\cite{carrero2026ifac}. Their dimensions
specify how valid instances satisfying \(\Phi_p(x)\) are systematically
varied during test generation and should therefore be interpreted as
experimental generation dimensions rather than as additional semantic
conditions of \(P_p\). In particular, they define how the generator explores
the valid input space without redefining the behavioural expectation
\(B_p(x,D(x))\).

The generation spaces, their experimental strata, admissibility conditions,
and partial-oracle expectations were fixed before executing the PBT campaign.
They were not tuned or redefined after observing SUT predictions, conformance
rates, or counterexamples. This separation ensures that the frozen SUT is
evaluated against predefined behavioural expectations rather than expectations
derived from the observed failures.

For each generated valid test case \(x=G_p(\theta)\), the SUT output \(D(x)\)
was evaluated through the corresponding property-specific partial oracle
\(O_p\), following the property-checking procedure defined in
Section~\ref{sec:property_checking}. A test case was recorded as conforming
when \(O_p(x,D(x))=\mathsf{pass}\); otherwise, it was retained as a
counterexample in \(\mathcal{CE}_p\). For each counterexample, the property,
target variable or sensor configuration, and generation parameters were
preserved to support subsequent behavioural characterization.

The resulting conformance and counterexample evidence provides the basis for
the behavioural evaluation of the SUT. The following subsection specifies how
the property-specific generation spaces were systematically exercised during
the experimental campaign.

\subsection{PBT Campaign Configuration}
\label{sec:pbt_campaign}

The PBT campaign followed a stratified test-generation strategy. For each
property \(P_p\), selected dimensions of the generation space \(\Theta_p\),
specified in Table~\ref{tab:pbt_generation_config}, were systematically
combined into predefined experimental configurations, hereafter referred to
as \emph{strata}. Each stratum therefore represents a specific
property-generation configuration applied to a set of admissible base
contexts. The strata were fixed before campaign execution; pseudo-random
sampling was used only to select the base contexts within each stratum.

The total number of property checks was determined by the resulting strata
and the number of admissible base contexts exercised within each stratum,
rather than by an arbitrary global sample size.
Table~\ref{tab:pbt_campaign_design} summarizes the stratified design and its
realized campaign size.

\begin{table}[!t]
\centering
\caption{Stratified design and realized size of the PBT campaign.}
\label{tab:pbt_campaign_design}
\scriptsize
\setlength{\tabcolsep}{2pt}
\renewcommand{\arraystretch}{1.15}

\begin{tabular}{|c|p{4.0cm}|c|c|c|c|}
\hline
\textbf{Property} &
\centering\arraybackslash\textbf{Stratum construction} &
\textbf{Strata} &
\textbf{Cases/stratum} &
\textbf{Planned} &
\textbf{Executed} \\
\hline

MV &
\(9\) sensors &
9 & 500 & 4,500 & 4,500 \\
\hline

SS &
\(8\) sensors \(\times\) \(5\) durations &
40 & up to 100 & 4,000 & 2,157 \\
\hline

N &
\((6\) sensors \(\times\) \(2\) directions \(+\)
\(3\) sensors \(\times\) \(1\) direction\()\)
\(\times\) \(5\) magnitudes &
75 & 80 & 6,000 & 6,000 \\
\hline

ORV &
\(9\) sensors \(\times\) \(2\) boundaries
\(\times\) \(4\) \(\alpha\) levels &
72 & 100 & 7,200 & 7,200 \\
\hline

CA &
\(3\) radiation \(+\) \(9\) CO$_2$--ventilation \(+\)
\(3\) temperature configurations &
15 & 100 & 1,500 & 1,500 \\
\hline

CD &
\(4\) pairs \(\times\) \(2\) reference conditions &
8 & 200 & 1,600 & 1,600 \\
\hline

PSF &
\(9\) variables, \(4\)--\(8\) missing &
381 & 10 & 3,810 & 3,810 \\
\hline

\multicolumn{4}{|r|}{\textbf{Total}} &
\textbf{28,610} &
\textbf{26,767} \\
\hline
\end{tabular}
\end{table}


SS requires specific consideration because its property checks are defined
over temporal windows and the number of admissible non-overlapping windows
therefore constrains the realizable campaign size. The nominal design
comprised eight sensors and five durations, resulting in 40 sensor--duration
strata with a target of 100 cases per stratum, i.e., 4,000 planned checks.
To avoid generating overlapping instances from the same temporal segment,
only non-overlapping admissible windows were considered. Accordingly, for
sensor \(k\) and duration \(r\), the realized number of cases was
\(n_{k,r}=\min(100,N_{k,r})\), where \(N_{k,r}\) denotes the number of
non-overlapping admissible windows available for that stratum. No sampling
with replacement was used and no applicability or admissibility condition
was relaxed to reach the nominal quota. As a result, all 40 predefined SS
strata were retained, but the number of executed SS checks was reduced from
the 4,000 initially planned to 2,157.

The PSF stratum count also requires clarification. As specified in
Table~\ref{tab:pbt_generation_config}, PSF involves between four and eight
missing variables among the nine monitored variables. Each distinct subset
of affected variables was considered a separate experimental stratum.
Exhaustive coverage therefore yields
\(\binom{9}{4}+\binom{9}{5}+\binom{9}{6}+\binom{9}{7}+
\binom{9}{8}=381\) strata. With ten admissible base contexts per stratum,
this results in 3,810 PSF property checks.


Within each predefined stratum, test cases were generated from admissible
base contexts satisfying the conditions of the corresponding property.
Base contexts were sampled without replacement using a fixed random seed,
making their selection reproducible while preserving the predefined
stratified design.

The campaign comprised 28,610 planned property checks. The only reduction
occurred for SS because of the non-overlapping-window constraint described
above, resulting in 26,767 executed checks. This value represents the
realized size of the stratified PBT campaign rather than a random sample
size drawn from the empirical dataset. Each generated case retained its
property identifier, stratum, generation parameters, selected base context,
SUT output, partial-oracle outcome, and counterexample status, providing
end-to-end traceability from test generation to property evaluation.

{\color{red}The PBT experimental layer was implemented in Python and executed against
the previously trained LightGBM-based SUT without retraining or modifying
the model during the campaign. The implementation, experimental
configuration, and scripts for test generation, property checking, campaign
execution, and result aggregation will be made publicly available upon
acceptance. These artifacts provide the software and configuration required
to reproduce the PBT workflow and inspect the generated evidence. Full
regeneration of the campaign from the original observations additionally
requires access to the underlying greenhouse dataset, which is not
distributed with these artifacts.
}
\section{Results}
\label{sec:results}

This section presents the empirical results of the PBT study, from the
conventional predictive baseline of the SUT to the property-specific
behavioural evidence revealed by the executed testing campaign.

\subsection{Conventional ML Evaluation of the SUT}
\label{sec:ml_baseline}

Before applying the proposed PBT framework, the predictive performance of
the ML-based SUT was established through conventional temporal validation.
The development data were chronologically divided into training and
validation folds, preserving temporal ordering rather than randomly
partitioning observations. Performance on the held-out temporal validation
fold was assessed using standard predictive metrics for anomaly detection. This evaluation provides the conventional
predictive baseline against which the complementary behavioural evidence
obtained through property-based testing can subsequently be interpreted.

Table~\ref{tab:model_performance} reports the performance of the final
LightGBM-based pipeline on the temporal validation fold. For the binary
detector, precision, recall, and F1-score refer to the anomaly class, which
constitutes the operationally relevant minority class. The operating
threshold was set to 0.80, selected to maximize anomaly F1-score while
maintaining recall above 0.97. For the semantic stage, the reported accuracy
corresponds to the final multiclass anomaly-type decision over all injected
anomalous validation points; missed binary detections are therefore counted
as ``no type'' errors.

\begin{table}[!t]
\centering
\caption{LightGBM supervised pipeline performance on the temporal validation fold.}
\label{tab:model_performance}
\renewcommand{\arraystretch}{1.2}
\resizebox{\linewidth}{!}{%
\begin{tabular}{|l|c|c|c|c|c|}
\hline
\textbf{Stage} & \textbf{Rows / anomalies} & \textbf{Accuracy} &
\textbf{Precision} & \textbf{Recall} & \textbf{F1} \\
\hline
Binary detector & 221,099 rows; 28,792 anomalies &
99.20\% & 96.75\% & 97.15\% & 0.9695 \\
Multiclass anomaly-type decision & 28,792 anomalous points &
97.07\% & -- & -- & -- \\
\hline
\end{tabular}
}
\end{table}


Overall, the conventional evaluation shows strong predictive performance
for both anomaly detection and semantic classification. These results therefore establish a strong predictive baseline from which the additional behavioural evidence provided by property-based testing can be examined.

\subsection{RQ1 -- Property Operationalization}
\label{sec:results_rq1}

\noindent\textit{To what extent can domain knowledge about sensor anomalies
be operationalized as executable properties for testing the ML-based SUT?}

The seven domain-grounded anomaly behaviours were successfully
operationalized as executable properties covering the three evidence scopes
defined in Section~\ref{sec:behavioural_knowledge}: observation-level
(MV and ORV), temporal (SS and N), and multivariate (CA, CD, and PSF).
Across the 26,767 property checks realized by the campaign, every executed
test case satisfied the operational conditions \(\Phi_p\) of its
corresponding property and could therefore be submitted to the SUT and
evaluated by the associated property-specific partial oracle.

End-to-end traceability was preserved for every property check. The
implementation retained the property and generation configuration, generated
test case, SUT output, and partial-oracle outcome, allowing each evaluation
to be traced back to the conditions under which the corresponding test case
was generated. The domain-grounded specifications were therefore realized
as an executable workflow connecting property-specific generation,
operational-condition satisfaction, SUT execution, and partial-oracle
evaluation.

These results answer RQ1 by showing that the domain knowledge represented by
all seven anomaly behaviours, spanning observation-level, temporal, and
multivariate evidence, could be operationalized into executable and traceable
tests across all realized property checks. Whether the SUT conforms to the
behavioural expectations encoded by these properties is examined separately
in RQ2.

\subsection{RQ2 -- Property-Based Behavioural Evaluation}
\label{sec:results_rq2}

\noindent\textit{How does the ML-based SUT behave across the property-defined
input space exercised by each executable property?}

Across the complete PBT campaign, 24,231 of the 26,767 property checks were
evaluated as \(\mathsf{pass}\) by their corresponding property-specific
partial oracles, while 2,536 produced counterexamples. This corresponds to
an aggregate conformance rate of 90.53\%. However, this overall value masks
substantial differences across properties, as shown in
Table~\ref{tab:rq2_conformance}.

\begin{table}[!t]
\centering
\caption{Behavioural conformance of the SUT by executable property.}
\label{tab:rq2_conformance}
\scriptsize
\setlength{\tabcolsep}{4pt}
\renewcommand{\arraystretch}{1.15}
\begin{tabular}{|c|c|c|c|c|}
\hline
\textbf{Property} &
\textbf{Checks} &
\textbf{Pass} &
\textbf{Counterexamples} &
\textbf{Conformance} \\
\hline
MV  & 4,500 & 4,500 & 0     & 100.00\% \\
SS  & 2,157 & 1,435 & 722   & 66.53\% \\
N   & 6,000 & 5,470 & 530   & 91.17\% \\
ORV & 7,200 & 7,200 & 0     & 100.00\% \\
CA  & 1,500 & 1,490 & 10    & 99.33\% \\
CD  & 1,600 & 326   & 1,274 & 20.38\% \\
PSF & 3,810 & 3,810 & 0     & 100.00\% \\
\hline
\textbf{Total} &
\textbf{26,767} &
\textbf{24,231} &
\textbf{2,536} &
\textbf{90.53\%} \\
\hline
\end{tabular}
\end{table}

The property-level results reveal markedly different behavioural responses. MV, ORV, and PSF achieved complete conformance with their corresponding executable properties over all checks performed in the campaign, with no counterexamples observed. CA also showed near-complete conformance,
with only 10 counterexamples among 1,500 checks. N exhibited high but
incomplete conformance, with 530 counterexamples among 6,000 checks. In
contrast, SS showed substantially lower conformance: 722 of its 2,157 checks
produced counterexamples. The largest deviation from the expected property
behaviour occurred for CD, for which 1,274 of 1,600 checks resulted in
counterexamples, corresponding to a conformance rate of 20.38\%.

Given the conjunctive behavioural expectation \(B_p(x,D(x))\),
counterexamples were further distinguished as binary misses when the expected
anomaly decision was not produced, or semantic type mismatches when an
anomaly was detected but the expected semantic category was not assigned.
Figure~\ref{fig:rq2_failure_modes} reports the relative proportion of 
these two counterexample outcomes for each property exhibiting non-conformance; \(n\)
indicates the corresponding number of counterexamples. For SS, 98.75\% of
the counterexamples corresponded to binary misses, with a similarly dominant
pattern for N (94.53\%). CD exhibited the opposite behaviour: 98.74\% of its
counterexamples corresponded to semantic type mismatches rather than missed
anomaly decisions. CA presented both failure modes, although this result
should be interpreted in light of the small number of counterexamples
observed (\(n=10\)).

\begin{figure}[!t]
\centering
\includegraphics[width=\linewidth]
{figures/rq2_counterexample_failure_modes_by_property.pdf}
\caption{Counterexample outcomes by executable property.}
\label{fig:rq2_failure_modes}
\end{figure}

Overall, these results answer RQ2 by showing that the SUT behaviour is
property-dependent both in the degree of conformance and in the nature of
the observed counterexamples. Non-conformance for SS and N arises
predominantly from missed anomaly decisions, whereas CD counterexamples
arise predominantly from semantic-category mismatches. This
property-specific evidence complements the strong aggregate predictive
performance reported in Table~\ref{tab:model_performance} by revealing
behavioural differences that are not apparent from conventional performance
metrics alone. RQ3 further examines how this behavioural variation is distributed across
the predefined generation strata representing different sensor anomaly
dynamics.

\subsection{RQ3 -- Within-space Behavioural Characterization}
\label{sec:results_rq3}

\noindent\textit{How does SUT conformance vary across predefined dimensions
of the property-specific generation spaces?}

To answer RQ3, conformance was examined across the predefined dimensions of
the property-specific generation spaces \(\Theta_p\) specified in
Table~\ref{tab:pbt_generation_config}. The analysis preserves the
experimental strata established before campaign execution and summarized in
Table~\ref{tab:pbt_campaign_design}, rather than introducing post-hoc
partitions of the observed counterexamples. Accordingly, variation in
conformance can be related directly to the generation dimensions exercised
for each executable property. No additional thresholds, fitted models, or
post-hoc decision rules are introduced in this analysis.

The stratum-level analysis revealed markedly different patterns across the
seven executable properties. MV, ORV, and PSF maintained complete
conformance across all evaluated configurations, and therefore showed no
variation in conformance within their respective generation spaces. CA
likewise showed only isolated deviations across its contextual
configurations. In contrast, N, SS, and CD exhibited substantial variation
across their predefined strata. These three properties were therefore
examined in greater detail with respect to their property-specific generation
dimensions.

For N, the stratified design enables conformance to be examined across
perturbation magnitude for each sensor and applicable perturbation
direction. Each of the 75 N strata contains 80 property checks, providing
equal support across all sensor--direction--magnitude configurations. As
shown in Figure~\ref{fig:rq3_noise_boundary}, this comparison reveals that
the effect of perturbation magnitude is strongly sensor-dependent. The
largest variation across the evaluated magnitude levels occurs for DVEXT,
for which conformance remains at 22.50\% for the two lowest magnitudes and
then progressively increases to 32.50\%, 47.50\%, and 77.50\%. DCO2EXT
also exhibits a marked magnitude-dependent response in both perturbation
directions: reduced conformance is concentrated at the lowest levels,
whereas complete or near-complete conformance is reached at the higher
magnitudes. XGRINT follows a similar pattern for positive perturbations,
increasing from 71.25\% at the lowest magnitude to above 90\% from the
second level onwards and reaching complete conformance at the highest
level.

\begin{figure*}[!t]
\centering
\includegraphics[width=\textwidth]
{figures/rq3_noise_boundary_detail.pdf}
\caption{Observed conformance across executed perturbation-magnitude strata
for the N property.}
\label{fig:rq3_noise_boundary}
\end{figure*}

By contrast, the remaining sensors exhibit substantially more stable
behaviour across the evaluated magnitude range. DGREXT remains consistently
above 90\% conformance, while DTEXT, DHEXT, XCO2INT, XHINT, and XTINT
maintain high conformance over most of their configurations, with
comparatively limited variation between consecutive magnitude levels. For
sensors supporting perturbations in both directions, differences between
positive and negative perturbations are most apparent at the lower
magnitudes and become small once high conformance is reached.

Taken together, these profiles show that variation within
\(\Theta_{\mathrm{N}}\) is not governed by a common magnitude effect across
sensors. DCO2EXT and XGRINT exhibit reduced conformance at lower perturbation
magnitudes followed by consistently high conformance as magnitude increases.
DVEXT shows the strongest magnitude-dependent variation: although conformance
progressively increases, it remains below 80\% even at the highest evaluated
level. Conversely, the remaining sensors maintain predominantly high
conformance across the tested magnitude range. Thus, within the evaluated
N generation space, conformance varies primarily with sensor and perturbation
magnitude, while perturbation direction introduces secondary differences for
some sensors.

\begin{figure*}[!t]
\centering
\includegraphics[width=\textwidth]
{figures/rq3_ss_boundary_detail_heatmap.pdf}
\caption{Observed conformance across executed sensor-specific stuck-duration
strata for the SS property.}
\label{fig:rq3_ss_boundary}
\end{figure*}

For SS, behavioural variation was examined across the 40 predefined
sensor--duration strata. Unlike N, however, the number of executed checks
is not uniform across all configurations because the non-overlapping-window
constraint limits the number of admissible temporal instances available
for some sensor--duration combinations, as described in
Subsection~\ref{sec:pbt_campaign}. Figure~\ref{fig:rq3_ss_boundary}
therefore reports both the conformance rate and the corresponding number
of executed checks \(n\) for every stratum; the actual sensor-specific
stuck duration \(d\) is also shown within each cell.



The resulting profiles reveal pronounced sensor-dependent behaviour.
DC\-O2EXT maintains 100\% conformance across all five evaluated durations,
with the first three configurations supported by 100 checks and the two
longest durations by 71 and 31 checks, respectively. XCO2INT likewise
maintains complete conformance across all five duration levels. At the
opposite extreme, XGRINT exhibits 0\% conformance at every evaluated
duration. This behaviour is already observed with 100 checks at the
shortest duration and persists as the number of available non-overlapping
windows decreases at longer durations. DHEXT and XHINT define intermediate
but comparatively stable behavioural regions: DHEXT remains between
76\% and 82\% conformance across its five duration levels, whereas XHINT
ranges from 39\% to 50\%, with all ten corresponding strata supported by
100 checks.

The remaining sensors require a different interpretation because their
available support decreases sharply as duration increases. For DTEXT,
DVEXT, and XTINT, the number of executed checks falls from 31 or 100 at
the shorter durations to only one or a few cases at the longest levels.
Consequently, extreme conformance values observed in these sparsely
supported configurations cannot be interpreted in the same way as the
persistent patterns observed for sensors with substantially larger support.
For example, the 100\% conformance observed for DVEXT and XTINT at their
longest evaluated duration is based on a single executed check, as is the
0\% value for DTEXT. These configurations are retained in the analysis
because they correspond to predefined strata, but the associated support
limits the strength of any duration-dependent interpretation.

Taken together, the SS results show that conformance variation is primarily
sensor-dependent rather than governed by a common effect of stuck duration.
Complete conformance persists across the evaluated duration range for
DCO2EXT and XCO2INT, whereas XGRINT remains non-conformant throughout its
entire evaluated range. DHEXT and XHINT maintain comparatively stable
intermediate conformance despite changes in duration. No consistent
duration-dependent pattern is therefore observed across sensors within
\(\Theta_{\mathrm{SS}}\). Instead, the affected sensor constitutes the
dominant source of variation, while the unequal availability of admissible
windows limits the characterization of duration effects for DTEXT, DVEXT,
and XTINT.


\begin{figure}[!t]
\centering
\includegraphics[width=\linewidth]{figures/rq3_cd_boundary_detail_heatmap_4x2.pdf}
\caption{Observed CD conformance across executed variable-pair and
reference-condition strata.}
\label{fig:rq3_cd_boundary}
\end{figure}


For CD, the stratified design enables conformance to be examined across
the four variable pairs and the two reference conditions defined in
\(\Theta_{\mathrm{CD}}\): the reference variable \(v_{i,\ell}\) above
\(p_{75,\ell}\) or below \(p_{25,\ell}\). Each of the eight predefined
strata contains 200 property checks, providing equal empirical support
and allowing conformance to be compared directly across pair--condition
configurations. Figure~\ref{fig:rq3_cd_boundary} shows the resulting
conformance pattern.



The results reveal strongly pair-dependent responses to the reference
condition. The most pronounced contrast occurs for
\texttt{DGREXT}--\texttt{XGRINT}, for which conformance reaches 100.0\%
when the reference variable is above \(p_{75}\), but falls to 6.0\% when
it is below \(p_{25}\). For \texttt{DTEXT}--\texttt{XTINT}, the same
change in reference condition reduces conformance from 26.5\% to 0.0\%.
By contrast, \texttt{XTINT}--\texttt{XHINT} exhibits the opposite
asymmetry: conformance is only 2.5\% for the above-\(p_{75}\) condition
but increases to 28.0\% for the below-\(p_{25}\) condition.
\texttt{DHEXT}--\texttt{XHINT} constitutes a different case, with
0.0\% conformance under both reference conditions.

These results indicate that the behavioural variation within
\(\Theta_{\mathrm{CD}}\) cannot be characterized by a uniform effect of
the reference condition across correlated variable pairs. Instead, conformance depends jointly on the
specific pair and on whether the reference variable lies above
\(p_{75,\ell}\) or below \(p_{25,\ell}\), with the target variable placed
at the corresponding incompatible quartile according to the sign of the
pairwise correlation. Moving between these two predefined configurations
can therefore produce a substantial change in SUT behaviour, no change at
all, or a change in the opposite direction, depending on the variable pair.

Importantly, the aggregate CD conformance of 20.38\% reported in RQ2
conceals this variability. Across the eight predefined strata, conformance
ranges from 0.0\% to 100.0\%, including the 100.0\% versus 6.0\% contrast
observed for \texttt{DGREXT}--\texttt{XGRINT} and the 2.5\% versus 28.0\%
contrast observed for \texttt{XTINT}--\texttt{XHINT}. Consequently, the 1,274 CD counterexamples reported in RQ2 are not
uniformly distributed across the property-defined input space, but are
concentrated differently depending on the variable pair and reference
condition. This illustrates why counterexamples provide information beyond the
aggregate conformance rate: they identify specific configurations of
\(\Theta_{\mathrm{CD}}\) in which the expected property behaviour is not
preserved.

Taken together, these results answer RQ3 by showing that SUT conformance
varies systematically across predefined dimensions of the property-specific
generation spaces, but that the relevant dimensions differ across
properties. For N, variation is primarily associated with the target sensor
and perturbation magnitude, with perturbation direction introducing
secondary differences. For SS, the affected sensor constitutes the dominant
source of variation, while duration effects are less consistent and, for
some sensors, constrained by the availability of admissible windows. For CD,
conformance depends jointly on the variable pair and the reference condition.
These patterns are observed within the predefined and executed generation
spaces and should not be interpreted as universal thresholds for the
underlying physical process. By exposing within-property variation that is
hidden by aggregate predictive and property-level conformance metrics, the
PBT analysis provides a more detailed characterization of the observable
behaviour of the SUT.

\section{Discussion}
\label{sec:discussion}

This section interprets the empirical findings in relation to the proposed
PBT framework and the research questions. It discusses the complementary
role of behavioural conformance in ML evaluation, the implications for
testing ML-based anomaly detectors, the role and limitations of partial
oracles, the generalizability of the proposed approach, and the main threats
to validity.

\subsection{What Does PBT Add to ML Evaluation?}
\label{sec:discussion_pbt_ml}

The empirical study highlights a distinction between predictive performance
and behavioural conformance. Conventional ML evaluation characterizes how accurately a learned system
predicts over a predefined validation distribution, whereas the proposed PBT
framework examines whether domain-grounded behavioural expectations are
preserved over predefined property-specific generation spaces. RQ1 establishes that such expectations can be operationalized as executable
properties, while RQ2 evaluates their behavioural conformance and RQ3
characterizes how that conformance varies across predefined dimensions of the
corresponding property-specific generation spaces.

This distinction also changes the role of test-input generation. Rather
than using inputs only to estimate predictive performance, PBT uses
property-specific generation spaces to deliberately exercise domain-relevant
conditions under which a domain-grounded behavioural expectation is defined.
Generation thus becomes part of the testing strategy, enabling
specification-guided exploration beyond distribution-based evaluation.

Counterexamples add a further diagnostic dimension by linking behavioural
violations to the property configurations that produce them. As illustrated
by RQ3, this traceability makes it possible to move beyond aggregate
conformance and identify configurations within a generation space in which expected behaviour
is preserved or violated. PBT therefore complements conventional ML evaluation with behavioural
evidence that is both property-specific and configuration-specific.

This contribution also differs from using metamorphic relations or
capability-oriented behavioural tests in isolation. Metamorphic testing
typically evaluates relations between source and follow-up executions,
whereas capability-oriented behavioural testing organizes test cases around
expected capabilities. In the proposed PBT workflow, executable properties
instead integrate the \emph{WHEN} conditions, property-specific generation
space, generated case, observable SUT response, partial oracle, and
counterexample record within a single traceable testing process. This
end-to-end traceability supports the within-space behavioural
characterization examined in RQ3.

\subsection{Implications for Testing ML-Based Anomaly Detectors}
\label{sec:discussion_testing_implications}

The study suggests that testing ML-based anomaly detectors requires
properties that reflect the semantics and evidence requirements of the
behaviours being tested. Generic perturbations may exercise a model, but
they do not necessarily represent meaningful behavioural obligations.
Domain-grounded properties provide a more structured testing target by
specifying both the conditions under which a behaviour should be evaluated
and the expected response of the SUT.

A related implication concerns test-input generation. The differences in
conformance across properties observed in RQ2, together with the
within-space variation characterized in RQ3, show that a uniform generation
strategy would provide an incomplete view of the SUT.
The relevant dimensions of the generation space depend on the behaviour
represented by each property, including the target variable, temporal extent,
perturbation magnitude, contextual conditions, and relationships between
variables.
Test generation should therefore be designed around each property rather
than through a common perturbation strategy.

Finally, counterexamples are most useful when they remain traceable to the
property and generation configuration from which they originate. Such
traceability turns a behavioural violation into an actionable testing
artifact: it identifies configurations under which the expected behaviour is not preserved and provides a concrete basis for subsequent diagnosis,
test-suite refinement, or model improvement. For ML-based anomaly
detectors, retaining this connection between property, generated input,
SUT output, and oracle outcome is therefore as important as detecting the
violation itself.

\subsection{Ground Truth and Partial Oracles}
\label{sec:discussion_partial_oracles}

A central challenge in testing ML-based anomaly detectors is that complete
ground truth is rarely available for arbitrary operational inputs. The
proposed framework does not eliminate this oracle problem. Instead, RQ1
shows how domain knowledge can be operationalized into executable properties
whose behavioural expectations can be evaluated through property-specific
partial oracles under explicitly defined conditions. Their evaluation
through these partial oracles subsequently provides the basis for the
conformance and counterexample evidence examined in RQ2. Each partial oracle
is therefore intentionally local: it determines whether the SUT response
conforms to the behaviour specified by a property, rather than establishing
the correctness of every possible prediction.

For generated test cases, the expected behaviour is defined before the SUT
is executed. Each property specifies the conditions that a valid test case
must satisfy and the behaviour expected from the SUT when those conditions
hold. The expected outcome therefore derives from the property specification,
while controlled test generation constructs inputs satisfying its operational
conditions; neither is derived from the SUT prediction itself. This makes it
possible to evaluate selected behaviours without requiring complete
ground-truth labels for the entire operational input space.

A counterexample should consequently be interpreted as evidence that the
SUT violates a specified domain-grounded behavioural expectation under the
generated conditions, not as proof that the model is universally incorrect
for that input. The strength of the conclusion therefore depends on the
validity of the property and the domain knowledge from which it was derived.
Partial oracles reduce the practical impact of incomplete ground truth by
making selected expectations testable, but they do not replace complete
ground truth or remove the need to justify the behavioural specifications
on which those expectations depend.

\subsection{Generalizability}
\label{sec:discussion_generalizability}

The empirical findings of this study are specific to the evaluated SUT,
dataset, and property-specific generation spaces, and should not be assumed
to transfer unchanged to other models or operational settings. In
particular, the within-space conformance patterns identified in RQ3
characterize the predefined and executed generation spaces and should not
be interpreted as universal behavioural thresholds. What is intended to be
generalizable is the testing framework rather than these empirical patterns.
Its core workflow---deriving executable properties from domain knowledge,
defining property-specific generation spaces, evaluating SUT responses
through partial oracles, and retaining counterexamples---is not tied to a
particular learning algorithm or sensor domain.

Applying the framework elsewhere would nevertheless require
domain-specific instantiation. Relevant properties, generation dimensions,
operational conditions, and partial oracles must be defined according to the
behaviours and knowledge available in the target system. Generalizability
therefore lies in the testing methodology and its separation between
property specification, test generation, and behavioural evaluation, not
in reusing the concrete properties or empirical results of this study.


\subsection{Threats to Validity}
\label{sec:threats_validity}

\paragraph{Construct validity}
The main threat concerns the adequacy of the executable properties and the
behavioural expectations evaluated through their partial oracles. These
properties capture specific behavioural expectations under explicitly
defined conditions, but they do not represent the complete set of correct
behaviours of the SUT. Consequently, a passing check establishes conformance
only with respect to the corresponding property, while a counterexample
indicates violation of that particular expectation. The expected outcomes
are defined independently of the SUT predictions and were specified before
campaign execution, reducing the risk of circular evaluation or post-hoc
property construction. However, the validity of the resulting evidence still
depends on the adequacy of the prior empirical taxonomy, physical constraints,
and domain assumptions encoded by each property.

\paragraph{Internal validity}
The main threat concerns the definition and coverage of the
property-specific generation spaces. The campaign used predefined generation
configurations and a fixed random seed, sampled cases without replacement,
and evaluated all generated cases against the same frozen SUT without
retraining or adaptation. Test configurations and oracle outcomes were
retained to preserve traceability and reproducibility, and no property,
generator, or oracle was changed after observing the counterexamples.
Nevertheless, the explored strata and parameter ranges determine which
configurations of each generation space are evaluated. For SS, coverage is
additionally constrained by the availability of non-overlapping admissible
temporal windows, resulting in unequal numbers of executed cases across some
sensor--duration strata. The behavioural variation identified in RQ3 should
therefore be interpreted within the executed configurations rather than as
an exhaustive characterization of the complete SUT input space.

\paragraph{External validity}
The empirical scope is limited to a single ML-based anomaly-processing
pipeline and one greenhouse data source. Different learning algorithms,
datasets, sensor infrastructures, or application domains may exhibit
different conformance patterns. The empirical findings should therefore not
be generalized beyond the evaluated setting without further replication.
As discussed in Section~\ref{sec:discussion_generalizability}, broader
applicability concerns the testing methodology rather than the specific
behavioural patterns observed in this case study. A related replication
threat concerns data availability: the public artifacts can document the
implementation, configuration, and generated evidence, but full regeneration
of the campaign requires access to the original greenhouse observations.


\section{Conclusions and Future Work}
\label{sec:conclusions}

This work presented a property-based testing framework for the behavioural
evaluation of ML-based anomaly detectors under incomplete ground truth.
The framework operationalizes domain knowledge as executable properties,
uses property-specific generation spaces to construct valid test cases, and
evaluates observable SUT behaviour through partial behavioural oracles.
The empirical study shows that this perspective complements conventional
predictive evaluation by revealing property-dependent behavioural conformance
and by characterizing how conformance varies across predefined dimensions of
the property-specific generation spaces. Counterexamples thereby provide
traceable, property- and configuration-specific testing evidence that is not
visible in aggregate predictive metrics alone.

Future work will evaluate the framework across alternative ML models,
datasets, and application settings, and investigate systematic support for
the derivation and specification of domain-grounded properties and generation
spaces. Further work will also study how counterexamples can guide test-suite
refinement and model improvement, while preserving the separation between
testing evidence and model adaptation.

\section*{Acknowledgements}
This work is the result of Project PID2024-157228OB-I00 of the 2024-2027 Spanish State Plan for Scientific and Technical Research and Innovation, co-financed by the Ministry of Science, Innovation, and Universities, the State Research Agency, and the European Regional Development Fund (ERDF).

\bibliographystyle{elsarticle-num}
\bibliography{references}

\end{document}
